\documentclass[12pt, a4paper, oneside]{book}

% Paquetes esenciales para matemáticas y formato
\usepackage[utf8]{inputenc}
\usepackage[T1]{fontenc}
% AGREGADO: 'es-noquoting' para evitar conflictos de < y > con las flechas de TikZ
\usepackage[spanish, es-tabla, es-noquoting]{babel} 
\usepackage{amsmath, amssymb, amsfonts, amsthm}
\usepackage{geometry}
\geometry{a4paper, top=3cm, bottom=3cm, left=3cm, right=3cm}
\usepackage{hyperref}
\usepackage{xcolor}
\usepackage{enumerate}
\usepackage{tikz}
\usepackage{caption}
\usetikzlibrary{babel, arrows.meta, positioning}

% Configuración de colores para los enlaces
\hypersetup{
	colorlinks=true,
	linkcolor=blue,
	filecolor=magenta,      
	urlcolor=cyan,
	citecolor=red,
}

% --- SISTEMA DE NUMERACIÓN INDEPENDIENTE POR SECCIÓN ---
\theoremstyle{definition}
\newtheorem{definicion}{Definición}[section]
\newtheorem{ejemplo}{Ejemplo}[section]
\newtheorem{ejercicio}{Ejercicio}[section]
\newtheorem{observacion}{Observación}[section]

\theoremstyle{plain}
\newtheorem{teorema}{Teorema}[section]
\newtheorem{proposicion}{Proposición}[section]
\newtheorem{corolario}{Corolario}[section]
\newtheorem{lema}{Lema}[section]

% Contadores de figuras y ecuaciones por sección
\numberwithin{figure}{section}
\numberwithin{equation}{section}

% Datos del documento
\title{
	\Huge \textbf{Notas de Seminario de Dinámica Topológica} \\
	\vspace{0.5cm}
	\Large Fundamentos y Aplicaciones
}
\author{
	\textbf{Matihus Molina Larios} \\
	\texttt{matihus.molina@unmsm.edu.pe} \\
	Facultad de Ciencias Matemáticas \\
	Universidad Nacional Mayor de San Marcos
}
\date{\today}

\begin{document}
	
	\frontmatter
	\maketitle
	
	\chapter*{Prefacio}
	\addcontentsline{toc}{chapter}{Prefacio}
	Estas notas de seminario sobre Dinámica Topológica han sido estructuradas como cimiento teórico para el estudio riguroso de los sistemas dinámicos. Los conceptos aquí desarrollados —tales como órbitas, conjuntos límite y recurrencia— sirven como base fundamental para avanzar hacia investigaciones más complejas. En particular, esta teoría proporciona las herramientas matemáticas necesarias para comprender la estructura subyacente en modelos avanzados, siendo un paso clave para futuras aplicaciones en la teoría ergódica.
	
	\tableofcontents
	
	\mainmatter
	
		\chapter{Conceptos Básicos en Sistemas Dinámicos}

\section{Espacios y Órbitas}

Consideremos un sistema dinámico discreto dado por un par $(M, f)$, donde $M$ puede ser un espacio topológico, una variedad diferenciable o un espacio medible, y $f: M \to M$ es una función (que será continua, diferenciable o medible, respectivamente). 

\begin{definicion}[Órbita futura y pasada]
	La \textbf{órbita futura} de $x \in M$ por $f$ se denota por $\mathcal{O}_f^+(x)$ y se define como el conjunto de iteraciones sucesivas:
	$$\mathcal{O}_f^+(x) = \{f^n(x) : n \in \mathbb{N} \cup \{0\}\}$$
	Si $f$ admite inversa (es decir, es biyectiva y $f^{-1}$ tiene las mismas propiedades de regularidad), definimos la \textbf{órbita pasada} de $x \in M$ como:
	$$\mathcal{O}_f^-(x) = \{f^{-n}(x) : n \in \mathbb{N} \cup \{0\}\}$$
	Y la \textbf{órbita total} como $\mathcal{O}_f(x) = \{f^n(x) : n \in \mathbb{Z}\}$.
\end{definicion}

\begin{ejemplo}
	Sea $f: \mathbb{R} \to \mathbb{R}$ dada por $f(x) = 2x$. 
	\begin{itemize}
		\item Para el punto $x = 0$, tenemos que $f(0) = 0$. Por lo tanto, su órbita futura y pasada trivializan en el origen: $\mathcal{O}_f^+(0) = \{0\}$.
		\item Para un punto $x_0 \neq 0$, por ejemplo $x_0 > 0$, las iteraciones sucesivas nos dan $f^n(x_0) = 2^n x_0$. Como se observa en el diagrama de telaraña de la Figura \ref{fig:orbita_2x}, al iterar la función, los puntos $x_1, x_2, \dots$ se alejan cada vez más, por lo que la órbita diverge hacia el infinito.
	\end{itemize}
\end{ejemplo}

\begin{figure}[htbp]
	\centering
	\begin{tikzpicture}[scale=2.5]
		% Cuadrícula de fondo
		\draw[lightgray, very thin] (-0.5,-0.5) grid (2.5,2.5);
		
		% Ejes
		\draw[thick, red, ->] (-0.8,0) -- (2.5,0) node[right, black] {$x$};
		\draw[thick, red, ->] (0,-0.8) -- (0,2.5) node[above, black] {$y$};
		
		% Recta y = x (identidad)
		\draw[thick, black] (-0.5,-0.5) -- (2.5,2.5) node[right] {$y=x$};
		
		% Función f(x) = 2x (en azul)
		\draw[thick, blue] (-0.25,-0.5) -- (1.25,2.5) node[right] {$f(x)=2x$};
		
		% Órbita (Diagrama de telaraña) partiendo de un x_0 pequeño
		\draw[dashed, black, thick] (0.2, 0) -- (0.2, 0.4); 
		\draw[dashed, black, thick] (0.2, 0.4) -- (0.4, 0.4); 
		\draw[dashed, black, thick] (0.4, 0.4) -- (0.4, 0.8); 
		\draw[dashed, black, thick] (0.4, 0.8) -- (0.8, 0.8); 
		\draw[dashed, black, thick] (0.8, 0.8) -- (0.8, 1.6); 
		\draw[dashed, black, thick] (0.8, 1.6) -- (1.6, 1.6); 
		
		% Puntos y etiquetas cuidadosamente espaciados
		\filldraw[black] (0,0) circle (0.6pt) node[below left] {$0$};
		\filldraw[black] (0.2,0) circle (0.6pt) node[below] {$x_0$};
		\filldraw[black] (0.4,0) circle (0.6pt) node[below] {$x_1$};
		\filldraw[black] (0.8,0) circle (0.6pt) node[below] {$x_2$};
	\end{tikzpicture}
	\caption{Gráfica de $f(x)=2x$ y diagrama de telaraña mostrando que la órbita $\mathcal{O}_f^+(x_0)$ diverge del origen.}
	\label{fig:orbita_2x}
\end{figure}

\section{Puntos Fijos y Periódicos}

\begin{definicion}[Puntos fijos y periódicos]
	El conjunto de los \textbf{puntos fijos} de $f$ se denota por $\text{Fix}(f)$ y se define por:
	$$\text{Fix}(f) = \{x \in M : f(x) = x\}$$
	El conjunto de los \textbf{puntos periódicos} de $f$ se denota por $\text{Per}(f)$ y se define por:
	$$\text{Per}(f) = \{x \in M : f^n(x) = x \text{ para algún } n \in \mathbb{N}\}$$
\end{definicion}

\begin{observacion}
	Se cumple que $\text{Fix}(f) \subset \text{Per}(f)$. En efecto, si tomamos $x \in \text{Fix}(f)$, por definición tenemos que $f(x) = x$. Esto significa que la condición de punto periódico $f^n(x) = x$ se satisface para el caso particular donde $n=1$. Por lo tanto, $x \in \text{Per}(f)$.
\end{observacion}

\begin{ejemplo}
	Considere la aplicación (conocida como el \textit{mapa tienda}) $f: [0,1] \to [0,1]$ definida por:
	$$ f(x) = \begin{cases} 
		2x & \text{si } 0 \le x < 1/2 \\ 
		2 - 2x & \text{si } 1/2 \le x \le 1 
	\end{cases} $$
	Halle explícitamente $\text{Fix}(f)$.
\end{ejemplo}

\begin{proof}[Solución]
	Para encontrar $\text{Fix}(f)$, resolvemos la ecuación $f(x)=x$ en ambos subintervalos:
	\begin{enumerate}
		\item Si $0 \le x < 1/2$, igualamos $2x = x \implies x = 0$.
		\item Si $1/2 \le x \le 1$, igualamos $2 - 2x = x \implies 3x = 2 \implies x = 2/3$.
	\end{enumerate}
	Por lo tanto, concluimos que $\text{Fix}(f) = \{0, 2/3\}$.
\end{proof}

\begin{figure}[htbp]
	\centering
	\begin{minipage}{0.31\textwidth}
		\centering
		\begin{tikzpicture}[scale=2.2]
			\draw[->] (-0.1,0) -- (1.1,0) node[right] {$x$};
			\draw[->] (0,-0.1) -- (0,1.1) node[above] {$f$};
			\draw[dashed, gray] (0,0) -- (1,1);
			\draw[thick, blue] (0,0) -- (0.5,1) -- (1,0);
			\filldraw[red] (0,0) circle (0.8pt) node[left] {$0$};
			\filldraw[red] (0.666,0.666) circle (0.8pt) node[right] {$\frac{2}{3}$};
		\end{tikzpicture}
		\caption*{Iteración $f(x)$}
	\end{minipage}\hfill
	\begin{minipage}{0.31\textwidth}
		\centering
		\begin{tikzpicture}[scale=2.2]
			\draw[->] (-0.1,0) -- (1.1,0) node[right] {$x$};
			\draw[->] (0,-0.1) -- (0,1.1) node[above] {$f^2$};
			\draw[dashed, gray] (0,0) -- (1,1);
			\draw[thick, blue] (0,0) -- (0.25,1) -- (0.5,0) -- (0.75,1) -- (1,0);
			\filldraw[red] (0,0) circle (0.8pt);
			\filldraw[red] (0.4,0.4) circle (0.8pt);
			\filldraw[red] (0.8,0.8) circle (0.8pt);
			\filldraw[red] (0.285,0.285) circle (0.8pt); 
			\filldraw[red] (0.666,0.666) circle (0.8pt);
		\end{tikzpicture}
		\caption*{Iteración $f^2(x)$}
	\end{minipage}\hfill
	\begin{minipage}{0.31\textwidth}
		\centering
		\begin{tikzpicture}[scale=2.2]
			\draw[->] (-0.1,0) -- (1.1,0) node[right] {$x$};
			\draw[->] (0,-0.1) -- (0,1.1) node[above] {$f^3$};
			\draw[dashed, gray] (0,0) -- (1,1);
			\draw[thick, blue] (0,0) -- (0.125,1) -- (0.25,0) -- (0.375,1) -- (0.5,0) -- (0.625,1) -- (0.75,0) -- (0.875,1) -- (1,0);
			% Solo marcando intersecciones referenciales
			\filldraw[red] (0,0) circle (0.8pt);
			\filldraw[red] (0.142,0.142) circle (0.8pt);
			\filldraw[red] (0.285,0.285) circle (0.8pt);
			\filldraw[red] (0.4,0.4) circle (0.8pt);
			\filldraw[red] (0.571,0.571) circle (0.8pt);
			\filldraw[red] (0.666,0.666) circle (0.8pt);
			\filldraw[red] (0.8,0.8) circle (0.8pt);
			\filldraw[red] (0.888,0.888) circle (0.8pt);
		\end{tikzpicture}
		\caption*{Iteración $f^3(x)$}
	\end{minipage}
	\caption{El mapa tienda y sus iteraciones sucesivas. Las intersecciones con $y=x$ (puntos rojos) revelan el crecimiento exponencial de los puntos periódicos.}
	\label{fig:mapa_tienda_iteraciones}
\end{figure}

\begin{ejercicio}
	Definamos $\text{Per}_n(f) = \{x \in M : f^n(x) = x\}$. Para el mapa tienda $f$ definido en el ejemplo anterior:
	\begin{enumerate}[a)]
		\item Halle explícitamente $\text{Per}_n(f)$.
		\item Muestre que $\overline{\text{Per}(f)} = [0,1]$.
		\item ¿Se cumple que $\text{Per}(f) = [0,1]$?
	\end{enumerate}
\end{ejercicio}

\begin{proof}[Solución]
	\textbf{a)} Podemos partir la función en 2 partes: $f_1(x) = 2x$ y $f_2(x) = 2 - 2x$. Las iteraciones de $f$ siguen un algoritmo constructivo donde, para calcular $f^n$, el intervalo se divide en $2^n$ partes iguales. En la Figura \ref{fig:algoritmo_f} observamos el árbol completo con el desglose hasta $f^3$, mostrando claramente la estructura de ramificación de este algoritmo.
	
	\begin{figure}[htbp]
		\centering
		% Usaremos resizebox por si la gráfica se sale ligeramente de los márgenes en algunas configuraciones, 
		% asegurando que el despliegue de las 8 ramas se vea impecable.
		\resizebox{\textwidth}{!}{
			\begin{tikzpicture}[
				box/.style={draw, thick, rounded corners, align=center, minimum width=2.5cm, minimum height=0.6cm},
				arrow/.style={-{Stealth}, thick, shorten >=2pt, shorten <=2pt}
				]
				% Nivel 1
				\node[box, text=blue, draw=blue] (f1) at (0, 2) {$f_1$};
				\node[box, text=blue, draw=blue] (f2) at (0, -2) {$f_2$};
				
				% Nivel 2
				\node[box] (f11) at (4, 3) {$f_1 \circ f_1$};
				\node[box] (f21) at (4, 1) {$f_2 \circ f_1$};
				\node[box] (f12) at (4, -1) {$f_1 \circ f_2$};
				\node[box] (f22) at (4, -3) {$f_2 \circ f_2$};
				
				% Nivel 3 (Completo y sin superposiciones)
				\node[box, text=red, draw=red] (f111) at (8.5, 3.5) {$f_1 \circ f_1 \circ f_1$};
				\node[box, text=red, draw=red] (f211) at (8.5, 2.5) {$f_2 \circ f_1 \circ f_1$};
				
				\node[box, text=red, draw=red] (f121) at (8.5, 1.5) {$f_1 \circ f_2 \circ f_1$};
				\node[box, text=red, draw=red] (f221) at (8.5, 0.5) {$f_2 \circ f_2 \circ f_1$};
				
				\node[box, text=red, draw=red] (f112) at (8.5, -0.5) {$f_1 \circ f_1 \circ f_2$};
				\node[box, text=red, draw=red] (f212) at (8.5, -1.5) {$f_2 \circ f_1 \circ f_2$};
				
				\node[box, text=red, draw=red] (f122) at (8.5, -2.5) {$f_1 \circ f_2 \circ f_2$};
				\node[box, text=red, draw=red] (f222) at (8.5, -3.5) {$f_2 \circ f_2 \circ f_2$};
				
				% Puntos suspensivos indicando la continuación del algoritmo
				\node[font=\bfseries\Huge] at (12, 0) {$\dots$};
				
				% Flechas Nivel 1 a 2
				\draw[arrow, blue] (f1.east) to[out=0, in=180] (f11.west);
				\draw[arrow, blue] (f1.east) to[out=0, in=180] (f21.west);
				\draw[arrow, blue] (f2.east) to[out=0, in=180] (f12.west);
				\draw[arrow, blue] (f2.east) to[out=0, in=180] (f22.west);
				
				% Flechas Nivel 2 a 3
				\draw[arrow, red] (f11.east) to[out=0, in=180] (f111.west);
				\draw[arrow, red] (f11.east) to[out=0, in=180] (f211.west);
				
				\draw[arrow, red] (f21.east) to[out=0, in=180] (f121.west);
				\draw[arrow, red] (f21.east) to[out=0, in=180] (f221.west);
				
				\draw[arrow, red] (f12.east) to[out=0, in=180] (f112.west);
				\draw[arrow, red] (f12.east) to[out=0, in=180] (f212.west);
				
				\draw[arrow, red] (f22.east) to[out=0, in=180] (f122.west);
				\draw[arrow, red] (f22.east) to[out=0, in=180] (f222.west);
				
				% Textos guía
				\node[font=\bfseries] at (4, 4.2) {$f^2$ (4 partes)};
				\node[font=\bfseries] at (8.5, 4.2) {$f^3$ (8 partes)};
			\end{tikzpicture}
		}
		\caption{Algoritmo de combinaciones ramificadas de $f_1$ y $f_2$ demostrando la expansión total hasta $f^3$.}
		\label{fig:algoritmo_f}
	\end{figure}
	
	Cada parte de $f^n$ se da por las combinaciones de $f_1$ y $f_2$ (es decir, composiciones de la forma $f_{j_n} \circ \dots \circ f_{j_1}$). Además, como tanto $f_1$ como $f_2$ son funciones lineales, la composición truncada en cada subintervalo es también una función lineal. Por ende, la solución de igualar esta composición truncada a la función identidad siempre nos dará a lo sumo un punto. Como consecuencia, $\text{Per}_n(f)$ tiene como máximo $2^n$ puntos periódicos distintos.
	
	\textbf{b)} Según el razonamiento algorítmico anterior, al dividir el intervalo $[0,1]$ en $2^n$ partes, sabemos que los puntos fijos de $f^n$ deben habitar cada uno de estos subintervalos. 
	Sean $x_0 \in (0,1)$ y $\epsilon > 0$ arbitrarios. Existe $N \in \mathbb{N}$ tal que $\frac{1}{2^N} < \epsilon$. Dividimos el intervalo en $2^N$ partes. Por ende, $x_0$ pertenece a algún subintervalo de longitud $\frac{1}{2^N}$. De esta manera, garantizamos que existe algún punto periódico $x$ cuya distancia a $x_0$ es como máximo $\frac{1}{2^N}$. 
	Luego, la distancia $|x - x_0| \le \frac{1}{2^N} < \epsilon$, por lo que $x \in V_\epsilon(x_0)$. Esto demuestra la densidad: $\overline{\text{Per}(f)} = [0,1]$.
	
	\textbf{c)} ¡No! Esto es incorrecto. Las composiciones truncadas de $f_1$ y $f_2$ tienen constantes racionales. Para algún $f^n_j$ en particular, tendrá esta regla de correspondencia: $f^n_j(x) = ax - b$, donde $a, b \in \mathbb{Q}$. El punto periódico es la solución a $ax - b = x$, de donde obtenemos $x = \frac{b}{a-1}$ (con $a \neq 1$). 
	Dado que $a,b \in \mathbb{Q}$, el punto periódico resultante será siempre un número racional. Como el intervalo $[0,1]$ contiene números irracionales, es imposible que $\text{Per}(f) = [0,1]$.
\end{proof}

\begin{ejemplo}
	Considere la rotación del círculo $S^1 \to S^1$, donde $S^1 = \{z \in \mathbb{C} : |z| = 1\}$. Definimos la aplicación $R_\theta(z) = z e^{2\pi i \theta}$, con un $\theta \in (0,1)$ fijo.
	Si evaluamos un ángulo racional, digamos $\theta = p/q \in \mathbb{Q}$, todos los puntos completarán ciclos exactos, por lo que es directo verificar que $\text{Per}(R_{p/q}) = S^1$. Esto se puede observar en la Figura \ref{fig:rotaciones}(a).
\end{ejemplo}

\begin{ejercicio}
	Para la aplicación $R_\theta: S^1 \to S^1$ definida anteriormente, demuestre que si $\theta \notin \mathbb{Q}$ (es irracional), entonces $\overline{\mathcal{O}^+_{R_\theta}(z)} = S^1$ para todo $z \in S^1$. (Ver visualización en la Figura \ref{fig:rotaciones}(b)).
\end{ejercicio}

\begin{figure}[htbp]
	\centering
	\begin{minipage}{0.45\textwidth}
		\centering
		\begin{tikzpicture}[scale=2]
			\draw[thick, ->] (-1.3,0) -- (1.3,0);
			\draw[thick, ->] (0,-1.3) -- (0,1.3);
			\draw[thick, blue] (0,0) circle (1);
			
			\filldraw[red] (1,0) circle (1.5pt) node[above right] {$z$};
			\filldraw[red] (-0.5,0.866) circle (1.5pt) node[above left] {$R_{1/3}(z)$};
			\filldraw[red] (-0.5,-0.866) circle (1.5pt) node[below left] {$R_{1/3}^2(z)$};
			
			% Flechas de rotacion
			\draw[->, dashed, red] (1.1,0) arc (0:120:1.1);
			\draw[->, dashed, red] (-0.55,0.952) arc (120:240:1.1);
			\draw[->, dashed, red] (-0.55,-0.952) arc (240:360:1.1);
		\end{tikzpicture}
		\caption*{(a) $\theta = 1/3 \in \mathbb{Q}$}
	\end{minipage}\hfill
	\begin{minipage}{0.45\textwidth}
		\centering
		\begin{tikzpicture}[scale=2]
			\draw[thick, ->] (-1.3,0) -- (1.3,0);
			\draw[thick, ->] (0,-1.3) -- (0,1.3);
			\draw[thick, blue] (0,0) circle (1);
			
			% Bucle para generar densidad irracional
			\foreach \i in {1,...,60} {
				\filldraw[red] ({cos(\i*137.5)}, {sin(\i*137.5)}) circle (0.8pt);
			}
			\filldraw[black] (1,0) circle (1.5pt) node[above right] {$z_0$};
		\end{tikzpicture}
		\caption*{(b) $\theta \notin \mathbb{Q}$ (Denso)}
	\end{minipage}
	\caption{Dinámica de la rotación en $S^1$. Si $\theta$ es racional, la órbita es finita y periódica. Si $\theta$ es irracional, la órbita envuelve el círculo denso.}
	\label{fig:rotaciones}
\end{figure}

\begin{proof}
	En lugar de pensar en $R_\theta: S^1 \to S^1$, pensaremos en la dinámica de una manera equivalente. Consideremos la función $f: \mathbb{R}/\mathbb{Z} \to \mathbb{R}/\mathbb{Z}$ (que identificamos con el intervalo $[0,1)$) donde $f([x]) = [x+\theta] \pmod 1$.
	
	Tomemos un $\epsilon > 0$. Existe $N \in \mathbb{N}$ tal que $1/N < \epsilon$. De esta manera, dividimos el intervalo $[0,1)$ en $N$ piezas de igual longitud: $[0, 1/N), [1/N, 2/N), \dots, [(N-1)/N, 1)$.
	
	Ahora tomamos los primeros $N+1$ puntos de la órbita de cualquier $x$ fijo: $x, f(x), f^2(x), \dots, f^N(x)$. Por el Principio del Palomar, como tenemos $N+1$ puntos y solo $N$ piezas, existen al menos 2 puntos que se ubican en la misma pieza. Es decir, existen $j < k$ con $0 \le j < k \le N$ tales que $f^j(x)$ y $f^k(x)$ están en el mismo subintervalo.
	
	La distancia entre estos dos puntos es $\delta = d(f^k(x), f^j(x)) \le 1/N < \epsilon$. Notemos matemáticamente que la diferencia es:
	$$f^k(x) - f^j(x) = [x+k\theta] - [x+j\theta] = [(k-j)\theta] \pmod 1$$
	Sea $m = k-j$. La aplicación $f^m$ avanza de manera pequeñísima una distancia de $\delta$. Y como $\delta$ es menor a $\epsilon$ (la longitud de cualquier vecindad de tamaño $1/N$), entonces es imposible que la iteración salte alguna pieza. Finalmente, la órbita acabará cayendo en algún momento en cualquier pieza o vecindad a la que queramos llegar.
	
	Para regresar al círculo unitario, sea $T: S^1 \to \mathbb{R}/\mathbb{Z}$ dada por $T(e^{2\pi i x}) = [x]$. Se tiene la relación $T \circ R_\theta = f \circ T$. 
	Por ende, sea $z_0 \in S^1$ un punto fijo y un abierto $U \subset S^1$. Su imagen $T(U)$ es un abierto en $\mathbb{R}/\mathbb{Z}$. Por la densidad demostrada arriba, existe algún $r \in \mathbb{N}$ tal que $f^{mr}(T(z_0)) \in T(U)$. 
	Como $f^{mr} \circ T = T \circ R_\theta^{mr}$, deducimos que $T(R_\theta^{mr}(z_0)) \in T(U)$, lo que implica que $R_\theta^{mr}(z_0) \in U$. Así, la órbita intersecta cualquier abierto, demostrando que $\overline{\mathcal{O}^+_{R_\theta}(z)} = S^1$.
\end{proof}

\section{Conjuntos Límite}

\begin{definicion}[Conjuntos Límite $\omega$ y $\alpha$]
	Sea $f: M \to M$ una aplicación continua y $M$ un espacio topológico de Hausdorff y localmente compacto. El \textbf{conjunto $\omega$-límite} de $x$ por $f$, se denota por $\omega_f(x)$ y se define por:
	$$\omega_f(x) = \{y \in M : \exists \text{ sucesión } (n_k) \subseteq \mathbb{Z}^+ / f^{n_k}(x) \xrightarrow{k \to \infty} y\}$$
	Si $f$ es invertible, podemos definir el \textbf{conjunto $\alpha$-límite} de $x$ por $f$ como:
	$$\alpha_f(x) = \{y \in M : \exists \text{ sucesión } (n_k) \subseteq \mathbb{Z}^+ / f^{-n_k}(x) \xrightarrow{k \to \infty} y\}$$
\end{definicion}

\begin{ejemplo}
	Analicemos los conjuntos límite para algunas funciones:
	\begin{itemize}
		\item \textbf{Para $f(x)=2x$}: 
		Si $x \neq 0$, la órbita diverge a infinito, por lo que $\omega_f(x) = \emptyset$.
		Si $x = 0$, la órbita es fija, por lo que $\omega_f(0) = \{0\} = \alpha_f(0)$.
		Para el límite pasado con $x \neq 0$, $f^{-n}(x) = x/2^n \to 0$, por ende $\alpha_f(x) = \{0\}$.
		
		\item \textbf{Para la rotación $R_{1/3}$}:
		Sabemos que $\mathcal{O}_{R_{1/3}}^+(z) = \{z, R_{1/3}(z), R_{1/3}^2(z)\}$.
		Como las iteraciones son periódicas ($R_{1/3}^{3k}(z) = z$, $R_{1/3}^{3k+1}(z) = R_{1/3}(z)$, etc.), concluimos que $\omega_{R_{1/3}}(z) = \mathcal{O}_{R_{1/3}}^+(z)$ y por invertibilidad $\alpha_{R_{1/3}}(z) = \mathcal{O}_{R_{1/3}}^+(z)$.
	\end{itemize}
\end{ejemplo}

\begin{ejercicio}
	Si $p \in \text{Per}(f)$ con periodo $k \in \mathbb{Z}^+$. Demuestre que $\omega_f(p) = \mathcal{O}_f^+(p)$.
\end{ejercicio}
\begin{proof}[Solución]
	Como $f^k(p) = p$ y $f^i(p) \neq p$ para todo $1 \le i < k$, tenemos que la órbita es finita: $\mathcal{O}_f^+(p) = \{p, f(p), \dots, f^{k-1}(p)\}$. Es claro que $\mathcal{O}_f^+(p) \subseteq \omega_f(p)$.
	Ahora, si $x \in \omega_f(p)$, existe una sucesión $n_j \to \infty$ tal que $f^{n_j}(p) \to x$. Como la órbita $\mathcal{O}_f^+(p)$ es un conjunto discreto y finito, la única forma de que la sucesión converja a $x$ es que eventualmente sea constante e igual a $x$. Luego, $x \in \mathcal{O}_f^+(p)$.
\end{proof}

\begin{definicion}[Conjunto Invariante]
	Decimos que $A \subseteq M$ es invariante (para adelante) por $f: M \to M$ si $f(A) \subseteq A$.
\end{definicion}

\begin{ejemplo}
	Los conjuntos $\text{Fix}(f)$ y $\text{Per}(f)$ son invariantes.
	\begin{itemize}
		\item Para $\text{Fix}(f)$: Sea $x \in \text{Fix}(f) \implies f(x) = x \in \text{Fix}(f)$. Por ende $f(\text{Fix}(f)) \subseteq \text{Fix}(f)$. Además $x \in f(\text{Fix}(f))$ porque $\exists y \in \text{Fix}(f)$ tal que $f(y)=x$, de hecho $y=x$. Así, $f(\text{Fix}(f)) = \text{Fix}(f)$.
		\item Para $\text{Per}(f)$: Sea $x \in \text{Per}(f) \implies \exists y \in \text{Per}(f)$ tal que $f(x)=y$. Como $x$ es periódico, $\exists k \in \mathbb{Z}^+$ tal que $f^k(x)=x$. Evaluando: $f^k(y) = f^k(f(x)) = f(f(f^{k-1}(x))) = f(x) = y$. Entonces $y \in \text{Per}(f)$.
	\end{itemize}
\end{ejemplo}

\begin{proposicion}
	Si $M$ es un espacio topológico de Hausdorff y localmente compacto, entonces para todo $x \in M$, el conjunto $\omega_f(x)$ es cerrado e invariante por $f$ (es decir, $f(\omega_f(x)) = \omega_f(x)$).
\end{proposicion}

\begin{proof}
	Primero, demostramos que $\omega_f(x)$ es cerrado. Por definición, el conjunto puede escribirse equivalentemente como:
	$$\omega_f(x) = \bigcap_{k \ge 1} \overline{\{f^n(x) : n \ge k\}}$$
	Al ser la intersección arbitraria de conjuntos cerrados, $\omega_f(x)$ es cerrado (AF1).
	
	Para probar la invarianza (AF2), sea $y \in \omega_f(x)$. Existe una sucesión $n_k \to \infty$ tal que $f^{n_k}(x) \to y$. Por la continuidad de $f$, deducimos que $f(f^{n_k}(x)) \to f(y)$, lo cual se reescribe como $f^{n_k+1}(x) \to f(y)$. Puesto que la sucesión de índices $n_k+1$ también tiende a infinito, concluimos que $f(y) \in \omega_f(x)$. 
\end{proof}

\begin{definicion}[Conjunto Límite Global]
	Definimos las uniones de todos los conjuntos límite como:
	$$L_+(f) = \overline{\bigcup_{x \in M} \omega_f(x)}, \quad L_-(f) = \overline{\bigcup_{x \in M} \alpha_f(x)}$$
	El conjunto límite global de $f$ se denota por $L(f)$ y se define por $L(f) = L_+(f) \cup L_-(f)$.
\end{definicion}

\begin{observacion}
	Se cumple la cadena de inclusiones: $\text{Fix}(f) \subseteq \text{Per}(f) \subseteq L(f)$.
	\textbf{Demostración:} Sea $x \in \text{Per}(f)$ con periodo $m \in \mathbb{Z}^+$. Como demostramos anteriormente, $\omega_f(x) = \mathcal{O}_f^+(x) = \{x, f(x), \dots, f^{m-1}(x)\}$. Dado que $x \in \omega_f(x)$, entonces $x \in \bigcup_{y \in M} \omega_f(y) \subseteq L_+(f) \subseteq L(f)$.
\end{observacion}

\begin{observacion}
	Las inclusiones anteriores pueden ser estrictas: $\text{Fix}(f) \subsetneq \text{Per}(f) \subsetneq L(f)$.
	Por ejemplo, para la rotación irracional $f = R_{\sqrt{2}/2}$, no existen puntos fijos ni periódicos ($\text{Fix}(f) = \emptyset, \text{Per}(f) = \emptyset$). Sin embargo, como $\overline{\mathcal{O}_{R_{\sqrt{2}/2}}^+(z)} = S^1$, tenemos que $\omega_{R_{\sqrt{2}/2}}(z) = S^1$, lo que implica que $L(f) = S^1$.
\end{observacion}

\begin{ejercicio}
	Demuestre que $L(f)$ es invariante para adelante por $f$ (asumiendo $f \in C^0$).
\end{ejercicio}
\begin{proof}[Solución]
	Sea $y \in f(L(f))$. Sin pérdida de generalidad (trabajando con $L_+$), esto implica que $y = f(x)$ donde $x \in L_+(f) = \overline{\bigcup_{z \in M} \omega_f(z)}$.
	Entonces $f(x) \in f\left(\overline{\bigcup_{z \in M} \omega_f(z)}\right)$. Por continuidad, $f(\overline{A}) \subseteq \overline{f(A)}$, por lo que:
	$$f(x) \in \overline{f\left(\bigcup_{z \in M} \omega_f(z)\right)} \subseteq \overline{\bigcup_{z \in M} f(\omega_f(z))}$$
	Como $\omega_f(z)$ es invariante, $f(\omega_f(z)) \subseteq \omega_f(z)$, lo que nos da:
	$$\overline{\bigcup_{z \in M} f(\omega_f(z))} \subseteq \overline{\bigcup_{z \in M} \omega_f(z)} = L_+(f) \subseteq L(f)$$
	Por ende, $f(L(f)) \subseteq L(f)$.
\end{proof}

\section{Puntos No Errantes}

\begin{definicion}[Conjunto No Errante]
	Sea $M$ un espacio topológico de Hausdorff. Decimos que $x \in M$ es un \textbf{punto no errante} si para cualquier vecindad abierta $U$ de $x$ y para todo $N \ge 1$, existe un $n \ge N$ tal que:
	$$f^n(U) \cap U \neq \emptyset$$
	El conjunto de todos los puntos no errantes de $f$ se denota por $\Omega(f)$. De igual forma, el conjunto de puntos errantes es su complemento, $\Omega(f)^c$.
\end{definicion}

\begin{proposicion}
	Se cumplen las siguientes propiedades para $\Omega(f)$:
	\begin{enumerate}[i)]
		\item $\Omega(f)$ es cerrado.
		\item $f(\Omega(f)) \subseteq \Omega(f)$ asumiendo que $f \in C^0$.
		\item Si $f$ es homeomorfismo $\implies f(\Omega(f)) = \Omega(f)$.
	\end{enumerate}
\end{proposicion}

\begin{proof}
	\textbf{i)} Demostraremos que el complemento $\Omega(f)^c$ es abierto en $M$. 
	Sea $x \in \Omega(f)^c$. Existe una vecindad abierta $U \subseteq M$ de $x$ tal que $f^n(U) \cap U = \emptyset$ para todo $n \in \mathbb{Z}^+$. Luego, cualquier elemento dentro de $U$ también cumple esta condición de punto errante, por ende $U \subseteq \Omega(f)^c$, probando que es abierto.
	
	\textbf{ii)} Si $y \in f(\Omega(f))$, existe $x \in \Omega(f)$ tal que $f(x) = y$. Sea $V_y$ una vecindad arbitraria de $y$. Como $f$ es continua, $f^{-1}(V_y)$ es una vecindad abierta de $x$.
	Dado que $x \in \Omega(f)$, existe $n \in \mathbb{Z}^+$ tal que $f^n(f^{-1}(V_y)) \cap f^{-1}(V_y) \neq \emptyset$.
	Esto implica que existe un elemento $z$ que pertenece a ambos conjuntos. Si $z \in f^{-1}(V_y)$, entonces $f(z) \in V_y$. Por otro lado, si $z \in f^n(f^{-1}(V_y))$, entonces $z = f^n(w)$ para algún $w \in f^{-1}(V_y)$. Evaluando $f$, tenemos $f(z) = f^{n+1}(w) = f^n(f(w))$. Como $f(w) \in V_y$, entonces $f(z) \in f^n(V_y)$.
	Por lo tanto, $f(z) \in f^n(V_y) \cap V_y \neq \emptyset$, demostrando que $y \in \Omega(f)$.
	
	\textbf{iii)} (Ejercicio) Si $f$ es homeomorfismo, aplicamos ii) a la función inversa: $f^{-1} \in C^0 \implies f^{-1}(\Omega(f)) \subseteq \Omega(f)$. Aplicando $f$ a ambos lados, obtenemos $\Omega(f) \subseteq f(\Omega(f))$. Combinando esto con la parte ii), concluimos que $f(\Omega(f)) = \Omega(f)$.
\end{proof}

\begin{proposicion}
	Se cumple la inclusión: $L(f) \subseteq \Omega(f)$.
\end{proposicion}

\begin{proof}
	Recordemos que $L(f) = L_+(f) \cup L_-(f)$. Trabajaremos con $y \in L_+(f) = \overline{\bigcup_{x \in M} \omega_f(x)}$ (el análisis para $L_-$ es análogo). Dividiremos la prueba en dos casos:
	
	\textbf{Caso 1:} $y \in \bigcup_{x \in M} \omega_f(x)$. Esto significa que $y \in \omega_f(x)$ para algún $x \in M$. Luego existe una sucesión $(n_k) \subseteq \mathbb{Z}^+$ tal que $f^{n_k}(x) \to y$. Sea $U$ una vecindad de $y$. Existen índices $m > n > 0$ tales que $f^n(x), f^m(x) \in U$.
	Así, $f^{m-n}(f^n(x)) = f^m(x) \in U$. Como $f^n(x) \in U$, entonces $f^{m-n}(U) \cap U \neq \emptyset$, por lo que $y \in \Omega(f)$.
	
	\textbf{Caso 2:} $y \in \partial\left(\bigcup_{x \in M} \omega_f(x)\right)$. Sea $U$ una vecindad de $y$. Entonces, $U \cap \left(\bigcup_{x \in M} \omega_f(x)\right) \neq \emptyset$.
	Esto implica que $U \cap \omega_f(x) \neq \emptyset$ para algún $x \in M$. Sea $z \in U \cap \omega_f(x)$. Como $z \in \omega_f(x)$, por el Caso 1 sabemos que $z \in \Omega(f)$. Dado que $U$ es una vecindad de $z$, entonces $f^k(U) \cap U \neq \emptyset$ para algún $k \in \mathbb{Z}^+$. Al ser $U$ una vecindad arbitraria de $y$, se concluye que $y \in \Omega(f)$.
\end{proof}

\begin{ejercicio}
	Halle $\Omega(f)$ para las siguientes aplicaciones:
	\begin{enumerate}[a)]
		\item $f(x)=2x$ en $\mathbb{R}$
		\item $f(z)=R_{1/3}(z)$ en $S^1$
		\item $f(z)=R_{\sqrt{3}/3}(z)$ en $S^1$
		\item $f(x) = \begin{cases} 2x & : 0 \le x < 1/2 \\ 2-2x & : 1/2 \le x \le 1 \end{cases}$ en $[0,1]$
	\end{enumerate}
\end{ejercicio}
\begin{proof}[Solución]
	\textbf{a)} Es claro que $0 \in \Omega(f)$. Supongamos por el absurdo que existe un $x_0 \neq 0$ tal que $x_0 \in \Omega(f)$. Sin pérdida de generalidad, asumamos $x_0 > 0$. Sea el abierto $U = (3x_0/4, 5x_0/4)$.
	Entonces, sus iteraciones son de la forma $f^n(U) = (2^n \cdot 3x_0/4, 2^n \cdot 5x_0/4)$. Para $n \in \mathbb{Z}^+$, se tiene que $2^n \ge 2$. Por lo tanto, el extremo inferior de $f^n(U)$ es $\frac{3 \cdot 2^n}{4} x_0 \ge \frac{6}{4} x_0 > \frac{5}{4} x_0$. 
	Luego, $f^n(U) \cap U = \emptyset$ para todo $n \in \mathbb{Z}^+$, lo cual contradice que $x_0 \in \Omega(f)$. Así, $\Omega(f) = \{0\}$.
	
	\textbf{b)} Sabemos que todos sus puntos son periódicos. Por ende, para todo $x \in S^1$ y cualquier vecindad $U_x$, se cumple que $f^q(U_x) \supseteq f^q(x) = x \in U_x$ para algún $q \in \mathbb{N}$ (en este caso $q=3$). Por tanto, $f^q(U_x) \cap U_x \neq \emptyset \implies \Omega(f) = S^1$.
	
	\textbf{c)} Como $\theta = \sqrt{3}/3 \notin \mathbb{Q}$, probamos antes que $S^1 = \overline{\mathcal{O}_f^+(x)}$ para todo $x \in S^1$. Es decir, para toda vecindad $U_x$ existe $f^n(x) \in U_x$ y además $f^n(x) \in f^n(U_x)$. Esto asegura la intersección, por lo que $\Omega(f) = S^1$. (También se pudieron resolver fácilmente sabiendo que $\omega_f(x) = S^1$ para todo $x \in S^1$ y aplicando el Teorema de inclusión $L(f) \subseteq \Omega(f)$).
	
	\textbf{d)} Sabemos que para el mapa tienda $\overline{\text{Per}(f)} = [0,1]$. Sea $x \in [0,1]$. Para cualquier vecindad $U_x$, existe por densidad un punto $y \in \text{Per}(f)$ tal que $y \in U_x$. Al ser periódico, existe $n \in \mathbb{N}$ tal que $f^n(y) = y$.
	Luego, $y \in f^n(U_x)$ y también $y \in U_x$, de modo que $U_x \cap f^n(U_x) \neq \emptyset$. Así, $\Omega(f) = [0,1]$.
\end{proof}
	% %%%%%%%%%%%%%%%%%%%%%%%%%%%%%%%%%%%%%%%%%%%%%%%%%%%%%%%%%%%%%%%%%%%%%%%%%
% %%% INICIO INPUT SEMANA 2 %%%
% (Corta desde aquí hasta la marca de "FIN INPUT SEMANA 2" 
% para pegarlo en "semana2.tex")
% %%%%%%%%%%%%%%%%%%%%%%%%%%%%%%%%%%%%%%%%%%%%%%%%%%%%%%%%%%%%%%%%%%%%%%%%%

\chapter{Semana 2: Rotaciones en el Círculo y Transitividad}

\section{Rotaciones en $S^1$}

Consideremos el cociente $\mathbb{R}/\mathbb{Z}$, donde $x \sim y \iff x - y \in \mathbb{Z}$. Este espacio puede ser identificado con el intervalo $[0,1)$ o con el círculo unidad $S^1 \subset \mathbb{C}$ mediante la biyección $[x] \mapsto e^{2\pi i x}$.

\begin{definicion}[Rotación en el Círculo]
	Dado un ángulo $\theta \in (0,1)$, definimos la \textbf{rotación de ángulo $\theta$} como la aplicación $R_\theta: \mathbb{R}/\mathbb{Z} \to \mathbb{R}/\mathbb{Z}$ dada por:
	$$R_\theta([x]) = [x + \theta]$$
	En la notación compleja de $S^1$, equivale a $R_\theta(z) = e^{2\pi i \theta} z$.
\end{definicion}

\begin{proposicion}
	Si $\theta \in \mathbb{Q}$, $\theta = p/q$ con $p, q \in \mathbb{Z}$, $q > 0$ y $(p,q)=1$ (fracción irreducible), entonces $R_\theta^q(z) = z, \forall z \in S^1$, y además $R_\theta^j(z) \neq z$ para $j=1,2,\dots,q-1$. Es decir, todos los puntos son periódicos de periodo exactamente $q$.
\end{proposicion}
\begin{proof}
	Tenemos $R_\theta^q(z) = e^{2\pi i \theta q} z = e^{2\pi i (p/q) q} z = e^{2\pi i p} z = z$.
	
	Supongamos que $\exists q_0 \in \{1,2,\dots,q-1\}$ tal que $R_\theta^{q_0}(z) = z$ para algún $z \in S^1$. Por el algoritmo de Euclides para la división: $q = q_0 \tau + r$, con $0 \le r < q_0$. 
	Luego, 
	$$z = R_\theta^q(z) = R_\theta^{q_0 \tau + r}(z) = R_\theta^r(R_\theta^{q_0 \tau}(z)) = R_\theta^r(z)$$
	Esto implica que $z = e^{2\pi i \theta r} z \implies 1 = e^{2\pi i \theta r} \implies \theta r = k \in \mathbb{Z}$. 
	Por tanto, $\theta = k/r = p/q \implies p r = k q$. Como $p$ y $q$ son primos entre sí ($(p,q)=1$), deducimos que $q$ debe dividir a $r$. Sin embargo, sabemos por el algoritmo de Euclides que $r < q_0 < q$, lo cual genera un absurdo, demostrando que $q$ es el menor entero positivo que cumple la condición.
\end{proof}

\begin{observacion}
	Si $\theta \in \mathbb{Q}$, la órbita de cualquier punto es finita y está dada por:
	$$\mathcal{O}_{R_\theta}(z) = \{R_\theta^n(z) : n \in \mathbb{Z}\} = \{z, R_\theta(z), \dots, R_\theta^{q-1}(z)\}$$
\end{observacion}

\begin{proposicion}
	Si $\theta \notin \mathbb{Q}$, la rotación $R_\theta$ no admite puntos periódicos.
\end{proposicion}

\begin{proof}
	Supongamos por contradicción que existe un punto periódico. Es decir, existen $[x] \in \mathbb{R}/\mathbb{Z}$ y $n \in \mathbb{N}$ tales que $R_\theta^n([x]) = [x]$. Por la definición de la rotación, esto implica que $[x + n\theta] = [x]$.
	Por lo tanto, la diferencia debe ser un número entero: $(x + n\theta) - x = k$ para algún $k \in \mathbb{Z}$. Esto implica que $n\theta = k \implies \theta = k/n$. Como $k, n \in \mathbb{Z}$ y $n \neq 0$, deducimos que $\theta \in \mathbb{Q}$, lo cual contradice la hipótesis de que $\theta$ es irracional.
\end{proof}

\begin{teorema}
	Si $\theta \notin \mathbb{Q}$, entonces para todo $z \in S^1$, la órbita futura es densa. Es decir: $\overline{\mathcal{O}_{R_\theta}^+(z)} = S^1$.
\end{teorema}

\begin{proof}
	Supongamos por el absurdo que $\overline{\mathcal{O}_{R_\theta}^+(z)} \neq S^1$. Denotemos $A = \overline{\mathcal{O}_{R_\theta}^+(z)}$. Como $A$ es cerrado propio, su complemento $A^c$ es abierto y no vacío, por lo que puede expresarse como la unión disjunta de sus componentes conexas, que son arcos abiertos en el círculo: $A^c = \bigcup (a_i, b_i)$.
	Dado que $A$ es un conjunto invariante bajo isometría (rotación), la distancia entre puntos se preserva. Es claro que $R_\theta(A) = A$. Tomando complementos, deducimos que $R_\theta(A^c) = A^c$.
	
	Como $R_\theta$ es un homeomorfismo, mapea componentes conexas en componentes conexas. Así, la imagen de un arco componente $(a, b) \subset A^c$ es otro arco $R_\theta(a, b) = (a+\theta, b+\theta) \subset A^c$. Al ser componentes conexas del mismo conjunto, solo hay dos opciones: o son exactamente el mismo arco, o son disjuntos.
	Si $R_\theta(a, b) = (a, b)$, entonces los extremos deben fijarse (bajo rotación completa, lo cual es imposible pues $\theta \notin \mathbb{Q}$). Por tanto, los arcos $R_\theta^n(a, b)$ son todos disjuntos para $n \in \mathbb{N}$.
	Como $R_\theta$ preserva la longitud, todos estos infinitos arcos disjuntos tienen la misma longitud $L = b-a > 0$. Sin embargo, la suma de las longitudes de infinitos arcos disjuntos daría infinito, lo cual es absurdo ya que todos están contenidos en $S^1$ que tiene longitud $2\pi$. Por lo tanto, $A$ debe ser todo $S^1$.
\end{proof}

\section{Conjugación Topológica}

\begin{definicion}[Conjugación Topológica]
	Sean $M, N$ espacios topológicos y $f: M \to M$, $g: N \to N$ aplicaciones continuas. Decimos que $f$ y $g$ son \textbf{topológicamente conjugadas} si existe un homeomorfismo $h: M \to N$ tal que $h \circ f = g \circ h$.
\end{definicion}

\begin{ejemplo}
	Considere la función tienda $f(x) = \begin{cases} 2x & 0 \le x < 1/2 \\ 2-2x & 1/2 \le x \le 1 \end{cases}$ en $[0,1]$ y la función logística $g(x) = 4x(1-x)$ en $[0,1]$. Muestre que $f$ y $g$ son topológicamente conjugadas y construya el homeomorfismo $h$.
\end{ejemplo}
\begin{proof}[Solución]
	Consideremos la función $h: [0,1] \to [0,1]$ dada por $h(x) = \sin^2\left(\frac{\pi x}{2}\right)$. 
	Primero verificamos que es biyectiva y continua: es continua al ser composición de seno, cuadrado y una función lineal. Su derivada es $h'(x) = \pi \sin(\frac{\pi x}{2}) \cos(\frac{\pi x}{2}) = \frac{\pi}{2} \sin(\pi x) > 0$ en $(0,1)$, por lo que es estrictamente creciente. Como $h(0)=0$ y $h(1)=1$, es una biyección continua de $[0,1]$ sobre sí mismo. 
	Por el teorema clásico de topología general, como $h$ es una biyección continua desde un espacio compacto ($[0,1]$) hacia un espacio de Hausdorff, esto garantiza que $h$ es un homeomorfismo.
	
	\begin{figure}[htbp]
		\centering
		\begin{minipage}{0.31\textwidth}
			\centering
			\begin{tikzpicture}[scale=2.5]
				\draw[->] (-0.1,0) -- (1.1,0) node[right] {$x$};
				\draw[->] (0,-0.1) -- (0,1.1) node[above] {$f(x)$};
				\draw[thick, blue] (0,0) -- (0.5,1) -- (1,0);
				\node[below left] at (0,0) {$0$};
				\node[below] at (1,0) {$1$};
			\end{tikzpicture}
			\caption*{Mapa Tienda $f(x)$}
		\end{minipage}\hfill
		\begin{minipage}{0.31\textwidth}
			\centering
			\begin{tikzpicture}[scale=2.5]
				\draw[->] (-0.1,0) -- (1.1,0) node[right] {$x$};
				\draw[->] (0,-0.1) -- (0,1.1) node[above] {$h(x)$};
				\draw[domain=0:1, smooth, variable=\x, purple, thick] plot ({\x}, {sin(\x*90)*sin(\x*90)});
				\node[below left] at (0,0) {$0$};
				\node[below] at (1,0) {$1$};
			\end{tikzpicture}
			\caption*{Conjugación $h(x)$}
		\end{minipage}\hfill
		\begin{minipage}{0.31\textwidth}
			\centering
			\begin{tikzpicture}[scale=2.5]
				\draw[->] (-0.1,0) -- (1.1,0) node[right] {$x$};
				\draw[->] (0,-0.1) -- (0,1.1) node[above] {$g(x)$};
				\draw[domain=0:1, smooth, variable=\x, red, thick] plot ({\x}, {4*\x*(1-\x)});
				\node[below left] at (0,0) {$0$};
				\node[below] at (1,0) {$1$};
			\end{tikzpicture}
			\caption*{Mapa Logístico $g(x)$}
		\end{minipage}
		\caption{Gráficas de la función tienda $f$, el homeomorfismo $h$ y la función logística $g$ que ilustran la conjugación topológica.}
		\label{fig:conj_tienda_logistica}
	\end{figure}
	
	Ahora verificamos la relación de conjugación $h(f(x)) = g(h(x))$:
	Por un lado, $g(h(x)) = 4\sin^2(\frac{\pi x}{2})\left[1 - \sin^2(\frac{\pi x}{2})\right] = 4\sin^2(\frac{\pi x}{2})\cos^2(\frac{\pi x}{2}) = \left(2\sin(\frac{\pi x}{2})\cos(\frac{\pi x}{2})\right)^2 = \sin^2(\pi x)$.
	
	Por otro lado, evaluamos $h(f(x))$ en ambos tramos de $f$:
	Si $x \in [0, 1/2)$, $h(f(x)) = \sin^2(\frac{\pi(2x)}{2}) = \sin^2(\pi x)$.
	Si $x \in [1/2, 1]$, $h(f(x)) = \sin^2(\frac{\pi(2-2x)}{2}) = \sin^2(\pi - \pi x)$. Por identidades trigonométricas, $\sin(\pi - \alpha) = \sin(\alpha)$, así que $\sin^2(\pi - \pi x) = \sin^2(\pi x)$.
	En ambos casos $h(f(x)) = \sin^2(\pi x) = g(h(x))$. Por lo tanto, son conjugadas.
\end{proof}

\begin{ejemplo}
	Considere $f: \mathbb{R} \to \mathbb{R}$ dada por $f(x) = \frac{1}{2}x$ y $g: \mathbb{R} \to \mathbb{R}$ dada por $g(x) = \frac{1}{3}x$. Muestre que $f$ y $g$ son topológicamente conjugadas, y construya el homeomorfismo $h: \mathbb{R} \to \mathbb{R}$ tal que $h \circ f = g \circ h$.
\end{ejemplo}
\begin{proof}[Solución]
	Necesitamos algún $h: \mathbb{R} \to \mathbb{R}$ tal que $h(x/2) = \frac{1}{3}h(x) \implies 3h(x/2) = h(x)$.
	Supongamos una función de la forma $h(x) = x^\alpha$. Sustituyendo, $3(x/2)^\alpha = x^\alpha \implies 3 = 2^\alpha \implies \alpha = \log_2 3$.
	
	Definimos entonces la función $h: \mathbb{R} \to \mathbb{R}$ considerando los signos:
	$$ h(x) = \begin{cases} 
		x^{\log_2 3} & : x \ge 0 \\ 
		-(-x)^{\log_2 3} & : x < 0 
	\end{cases} $$
	Claramente conjuga la dinámica. Además, al igual que sus funciones generadoras en sus respectivos intervalos (e.g. $[0, +\infty)$ y $(-\infty, 0)$), es claro que $h$ es biyectiva y continua.
\end{proof}

\begin{teorema}
	Sean $M, N$ espacios topológicos y $f: M \to M, g: N \to N$ aplicaciones continuas topológicamente conjugadas vía el homeomorfismo $h: M \to N$. Entonces:
	\begin{enumerate}[1)]
		\item $p \in M$ es punto periódico de $f$ de periodo "$n$" $\iff h(p)$ es un punto periódico de $g$ de periodo "$n$".
		\item $h(\mathcal{O}_f^+(x)) = \mathcal{O}_g^+(h(x)), \forall x \in M$ y $h(\mathcal{O}_f(x)) = \mathcal{O}_g(h(x))$ cuando $f$ y $g$ son homeomorfismos.
		\item $h(\omega_f(x)) = \omega_g(h(x)), \forall x \in M$ y $h(\alpha_f(x)) = \alpha_g(h(x))$ cuando $f$ y $g$ son homeomorfismos.
		\item $h(\Omega(f)) = \Omega(g)$.
	\end{enumerate}
\end{teorema}

\begin{proof}
	\textbf{1)} $[\Rightarrow]$ Sea $f^n(p)=p$ y $f^j(p) \neq p, \forall j \in \{1,\dots,n-1\}$. Tenemos: $h \circ f = g \circ h \implies h \circ f \circ f = g \circ h \circ f = g \circ g \circ h \implies h \circ f^2 = g^2 \circ h$. Generalizando, obtenemos que $h \circ f^n = g^n \circ h, \forall n \in \mathbb{Z}^+$.
	Luego: $g^n(h(p)) = h(f^n(p)) = h(p)$. 
	Supongamos que $\exists j \in \{1,\dots,n-1\}$ tal que $g^j(h(p)) = h(p) \implies h(f^j(p)) = h(p) \implies f^j(p) = p$, lo cual es absurdo ya que $n$ es el periodo mínimo de $p$. Así, "$n$" es el periodo de $h(p)$ respecto a $g$.
	
	$[\Leftarrow]$ Por hipótesis: $g^n(h(p)) = h(p)$ y $g^j(h(p)) \neq h(p), \forall j \in \{1,\dots,n-1\}$. Luego, $h(f^n(p)) = h(p) \implies f^n(p) = p$. Supongamos que $\exists j \in \{1,\dots,n-1\}$ tal que $f^j(p) = p \implies h(f^j(p)) = h(p) \implies g^j(h(p)) = h(p)$, un absurdo.
	
	\textbf{2)} Sea $y \in h(\mathcal{O}_f^+(x)) \iff \exists n \in \mathbb{N} / y = h(f^n(x)) \iff \exists n \in \mathbb{N} / y = g^n(h(x)) \iff y \in \mathcal{O}_g^+(h(x))$. Si $f$ y $g$ son homeomorfismos, la demostración para la órbita total es idéntica solo que usamos $\exists n \in \mathbb{Z} / y = h(f^n(x))$.
	
	\textbf{3)} Sea $y \in h(\omega_f(x)) \iff \exists (n_k) \subseteq \mathbb{Z}^+ \land z \in M$ tales que $f^{n_k}(x) \to z \land y = h(z) \iff$ [como $h$ es homeomorfismo] $\exists (n_k) \subseteq \mathbb{Z}^+ / h(f^{n_k}(x)) \to h(z) = y \iff g^{n_k}(h(x)) \to y \iff y \in \omega_g(h(x))$.
	
	\textbf{4)} Sea $y \in h(\Omega(f)) \implies \exists z \in \Omega(f) / h(z) = y$. Sea $V_y$ una vecindad de $y$, entonces $h^{-1}(V_y)$ es un conjunto abierto en $M$.
	Por definición de $\Omega(f)$, $\exists n \in \mathbb{Z}^+ / f^n(h^{-1}(V_y)) \cap h^{-1}(V_y) \neq \emptyset \implies \exists n \in \mathbb{Z}^+ / h(f^n(h^{-1}(V_y))) \cap h(h^{-1}(V_y)) \neq \emptyset \implies g^n(V_y) \cap V_y \neq \emptyset \implies y \in \Omega(g)$. La vuelta es completamente análoga.
\end{proof}

\begin{corolario}
	Sean $R_{r/s}$ y $R_{p/q}$ rotaciones racionales con $r,s,p,q \in \mathbb{Z} \land s,q > 0$ y fracciones irreducibles $(r,s)=1, (p,q)=1$. Entonces: 
	$$R_{r/s} \text{ y } R_{p/q} \text{ son topológicamente conjugadas} \implies s = q$$
\end{corolario}
\begin{proof}
	Sea $x \in S^1$. La cardinalidad de su órbita será exactamente el denominador del ángulo irreducible, es decir, $\text{card}(\mathcal{O}_{R_{r/s}}^+(x)) = s$. 
	Como $h$ es una función biyectiva, preserva la cantidad de elementos, lo que implica que $\text{card}(h(\mathcal{O}_{R_{r/s}}^+(x))) = s$.
	Por conjugación, esto debe ser igual a la cardinalidad de la órbita de $h(x)$ bajo la otra rotación: $\text{card}(\mathcal{O}_{R_{p/q}}^+(h(x))) = q$. Por consiguiente, se concluye que $s = q$.
\end{proof}

\begin{ejercicio}
	Mostrar que: $R_\theta$ y $R_\alpha$ son topológicamente conjugados $\iff \alpha \pm \theta \in \mathbb{Z}$.
\end{ejercicio}

\begin{ejercicio}
	Considere $f: S^1 \to S^1$, definida por $z \mapsto f(z) = z^2$.
	\begin{enumerate}[a)]
		\item Halle $\text{Per}_n(f) = \{z \in S^1 : f^n(z) = z\}$.
		\item Muestre que $\overline{\text{Per}(f)} = S^1$.
	\end{enumerate}
\end{ejercicio}
\begin{proof}[Solución]
	\textbf{a)} Para $n \in \mathbb{N} \land z \in \text{Per}_n(f) \implies f^n(z) = z^{2^n} = z \implies z^{2^n} - z = 0 \implies z(z^{2^n-1} - 1) = 0$.
	De aquí obtenemos $z=0 \lor z$ es una de las $2^n-1$ raíces de la unidad. Como estamos operando en el círculo unidad ($z \in S^1$), descartamos $z=0$. 
	Por tanto, $z = e^{2\pi i \frac{k}{2^n-1}} = \cos\left(2\pi \frac{k}{2^n-1}\right) + i\sin\left(2\pi \frac{k}{2^n-1}\right)$ para $0 \le k < 2^n-1$.
	
	\textbf{b)} Sabemos que las $N$ raíces de la unidad dividen al círculo uniformemente en $N$ partes de igual longitud ($L = 2\pi/N$). Sea un punto cualquiera $z = e^{2\pi i x} \in S^1$.
	Consideramos la biyección $h: [0,1) \to S^1, x \mapsto h(x) = e^{2\pi i x}$. 
	Sea $A$ un conjunto abierto que contiene a $z \implies h^{-1}(A)$ es un abierto que contiene a $x$. Luego existe un $\epsilon > 0 / V_\epsilon(x) \subseteq h^{-1}(A)$. 
	Sea $N \in \mathbb{N}$ un entero suficientemente grande tal que $1/N < \epsilon$ (conceptualmente este $N$ representará los valores $2^n-1$ para $n$ grandes). 
	
	Si $x = j/N$ para algún $j \in \{0, \dots, N-1\} \implies z = e^{2\pi i j/N} \implies z \in \text{Per}_N(f)$, por lo que la vecindad automáticamente lo contiene.
	Si $x \neq j/N, \forall j \in \{0, \dots, N-1\}$, partimos el intervalo $[0,1)$ en $N$ partes iguales. Entonces $x$ pertenece a algún subintervalo, lo que implica que su distancia a alguno de los extremos $j_0/N$ es estrictamente menor a la longitud de la partición: $|x - j_0/N| < 1/N < \epsilon$ para algún $j_0 \in \{0, \dots, N-1\}$.
	Esto asegura que $j_0/N \in V_\epsilon(x)$. Luego, al mapear de regreso, $y = h(j_0/N) \in A$, donde $y$ es evidentemente una raíz de la unidad $y \in \text{Per}_N(f)$. Esto demuestra que cualquier abierto interseca a los puntos periódicos, garantizando la densidad $\overline{\text{Per}(f)} = S^1$.
\end{proof}

\section{Sistemas Dinámicos Topológicamente Transitivos y Mixing}

\begin{definicion}[Topológicamente Transitivo]
	Sea $X$ un espacio topológico. Decimos que $f: X \to X$ es \textbf{topológicamente transitiva} si existe un $x \in X$ tal que su órbita futura es densa, es decir, $\overline{\mathcal{O}^+(x)} = X$.
\end{definicion}

\begin{ejemplo}
	$R_\theta: S^1 \to S^1$ donde $\theta \in \mathbb{R} \setminus \mathbb{Q}$ es topológicamente transitivo, pues $\overline{\mathcal{O}_{R_\theta}^+(z)} = S^1, \forall z \in S^1$.
\end{ejemplo}

\begin{ejemplo}
	$R_\theta: S^1 \to S^1$ donde $\theta \in \mathbb{Q}$ ($\theta = p/q$) implica que $\overline{\mathcal{O}_{R_\theta}^+(z)} = \{z, R_\theta(z), \dots, R_\theta^{q-1}(z)\} \subsetneq S^1$. Así, una rotación racional no es topológicamente transitiva.
\end{ejemplo}

\begin{teorema}
	Sea $X$ un espacio métrico separable (admite base numerable) y localmente compacto, y $f: X \to X$ continua. La función $f$ es topológicamente transitiva si y solo si para cualquier par de abiertos $U, V \subseteq X$ no vacíos, existe $n \in \mathbb{N}$ tal que $f^n(U) \cap V \neq \emptyset$.
\end{teorema}
\begin{proof}
	$(\Rightarrow)$ Supongamos que $f$ es topológicamente transitiva. Entonces existe $x \in X$ con órbita densa. Sean $U, V$ abiertos no vacíos. Por densidad, la órbita debe visitar ambos abiertos, es decir, existen enteros $j < k$ tales que $f^j(x) \in U$ y $f^k(x) \in V$.
	Aplicando $f^{k-j}$ a $f^j(x)$, obtenemos $f^{k-j}(f^j(x)) = f^k(x) \in V$. Como $f^j(x) \in U$, esto demuestra que $f^{k-j}(U) \cap V \neq \emptyset$, con $n = k-j > 0$.
	
	$(\Leftarrow)$ Asumimos la condición de intersección. Como $X$ es separable, admite una base numerable de abiertos $\{B_i\}_{i \in \mathbb{N}}$. Para cada $B_i$, el conjunto $U_i = \bigcup_{n \ge 0} f^{-n}(B_i)$ es la unión de preimágenes de un abierto, por ende es abierto. Afirmamos que cada $U_i$ es denso en $X$. En efecto, si tomamos cualquier abierto $W \subset X$, por hipótesis existe $n$ tal que $f^n(W) \cap B_i \neq \emptyset \implies W \cap f^{-n}(B_i) \neq \emptyset \implies W \cap U_i \neq \emptyset$, probando la densidad.
	Por el Teorema de Categoría de Baire, en un espacio localmente compacto de Hausdorff, la intersección numerable de abiertos densos sigue siendo un conjunto denso. Sea $x \in \bigcap U_i$. Esto implica que para todo básico $B_i$, $x \in \bigcup f^{-n}(B_i) \implies \exists n \text{ t.q. } f^n(x) \in B_i$. Así, la órbita de $x$ visita todo abierto básico, por lo que es densa y $f$ es topológicamente transitiva.
\end{proof}

\begin{ejercicio}
	Si $f: X \to X$ es topológicamente transitiva y $\varphi: X \to \mathbb{R}$ es continua tal que $\varphi(f(x)) = \varphi(x), \forall x \in X \implies \varphi$ es constante.
\end{ejercicio}
\begin{proof}[Solución]
	Supongamos que $\varphi$ no es constante $\implies \exists x, y \in X$ tales que $\varphi(x) \neq \varphi(y)$. 
	Sin pérdida de generalidad, asumamos que $\varphi(x) > \varphi(y) \implies \exists \alpha \in \mathbb{R}$ tal que $\varphi(x) > \alpha > \varphi(y)$.
	Definimos los conjuntos $U = \{z : \varphi(z) > \alpha\}$ y $V = \{z : \varphi(z) < \alpha\}$. Como $\varphi$ es continua, $U$ y $V$ son abiertos, y por la evaluación anterior, $\neq \emptyset$.
	Luego, como $f$ es topológicamente transitiva $\implies \exists k \in \mathbb{Z}^+ / f^k(U) \cap V \neq \emptyset$.
	Es decir, $\exists z \in V$ tal que $z = f^k(a)$ para algún $a \in U$. 
	Como se nos da que $\varphi(f^n(x)) = \varphi(x), \forall n \in \mathbb{N}$, evaluamos en $a$:
	$$ \varphi(z) = \varphi(f^k(a)) = \varphi(a) $$
	Dado que $z \in V \implies \varphi(z) < \alpha$. Pero como $a \in U \implies \varphi(a) > \alpha$. Por ende, obtenemos la contradicción $\varphi(z) > \alpha$ y $\varphi(z) < \alpha$ simultáneamente ($\rightarrow \leftarrow$).
	Por consiguiente, $\varphi$ debe ser constante.
\end{proof}

\begin{ejercicio}
	Sea $X$ un espacio separable y localmente compacto. Sea $f: X \to X$ continua. Mostrar que $f$ no es topológicamente transitiva $\iff \exists U, V \subset X$ abiertos t.q. $U \cap V = \emptyset$, $f(U) \subset U$ y $f(V) \subset V$.
\end{ejercicio}
\begin{proof}[Solución]
	$(\Leftarrow)$ Sean $U, V \subseteq X$ abiertos tales que $U \cap V = \emptyset$.
	Considerando $U' = \bigcup_{n \in \mathbb{N}} f^n(U)$ y $V' = \bigcup_{n \in \mathbb{N}} f^n(V)$, estos son conjuntos abiertos en $X$.
	Notemos que, como $f(U) \subseteq U$, se cumple por inducción que $f^k(U) \subseteq U$ para todo $k \in \mathbb{N}$. 
	De esto se deduce que:
	$$ f^k(U) \cap V \subseteq U \cap V = \emptyset \implies f^k(U) \cap V = \emptyset $$
	Al existir dos abiertos disjuntos cuya iteración nunca se interseca, se concluye que $f$ no es topológicamente transitiva.
	
	$(\Rightarrow)$ Sea $U_0, V_0 \subseteq X$ abiertos. Por la definición de no ser topológicamente transitiva, $\forall k \in \mathbb{Z}^+$ se cumple que $f^k(U_0) \cap V_0 = \emptyset$.
	Considerando $U = \bigcup_{n \in \mathbb{N}} f^n(U_0) \implies U$ es abierto.
	Además, notemos la invariancia: 
	$$f(U) \subseteq f\left(\bigcup_{n \in \mathbb{N}} f^n(U_0)\right) = \bigcup_{n \in \mathbb{N}} f^{n+1}(U_0) \subseteq U$$
	Sea $V = \text{int}(U^c) \implies$ como $X$ es localmente compacto, $V$ es abierto. Cumpliendo así la construcción de conjuntos invariantes y disjuntos requerida.
\end{proof}

\begin{definicion}[Topológicamente Mixing]
	Sea $X$ un espacio topológico y $f: X \to X$ continua. Decimos que $f$ es \textbf{topológicamente mixing} si para cualquier par de abiertos $U, V \subseteq X$ no vacíos, existe un $N \in \mathbb{N}$ tal que para todo $n \ge N$, $f^n(U) \cap V \neq \emptyset$.
\end{definicion}

\begin{observacion}
	Es claro que Topológicamente Mixing $\implies$ Topológicamente Transitivo, ya que la existencia de un $N$ para todos los $n \ge N$ implica fuertemente la existencia de al menos un $n$. El converso no siempre es cierto (por ejemplo, las rotaciones irracionales son transitivas pero no mixing, ya que los intervalos rotan rígidamente sin dispersarse por todo el espacio).
\end{observacion}

\begin{figure}[htbp]
	\centering
	\begin{tikzpicture}[scale=1.3]
		% Conjunto U original
		\draw[blue, thick, rounded corners=8pt] (0,0) -- (0.3,1.2) -- (1.5,1.5) -- (1.2,-0.2) -- cycle;
		\node[blue] at (0.7, 0.6) {$U$};
		
		% Flecha indicando iteración
		\draw[->, thick, gray, dashed] (1.8, 0.5) -- (2.8, 0.5) node[midway, above] {$f^n$};
		
		% f^n(U) expandido
		\draw[blue, thick, rounded corners=12pt] (3.2,0.2) -- (4,1.8) -- (8.5,1.5) -- (8,-0.8) -- cycle;
		\node[blue] at (5.5, 0.5) {$f^n(U)$};
		
		% Conjunto V (intersectado)
		\draw[red, thick, rounded corners=6pt] (7.5, 0.8) -- (8, 2) -- (9.5, 1.6) -- (9, 0.5) -- cycle;
		\node[red] at (8.6, 1.2) {$V$};
		\node[red, right] at (9.5, 1.2) {$\forall n \ge n_1$};
		
		% Conjunto W (intersectado)
		\draw[red, thick, rounded corners=6pt] (3.5, -0.6) -- (4.2, -1.2) -- (5.5, -0.5) -- (4.8, 0.2) -- cycle;
		\node[red] at (4.5, -0.4) {$W$};
		\node[red, right] at (4.5, -1) {$\forall n \ge n_2$};
		
	\end{tikzpicture}
	\caption{Dinámica de un sistema topológicamente mixing. La iteración $f^n(U)$ se expande lo suficiente como para intersecar de forma permanente cualquier otro abierto ($V$, $W$) a partir de cierto umbral temporal.}
	\label{fig:top_mixing_minimalista}
\end{figure}

\begin{ejemplo}
	Consideremos el mapa tienda $f: [0,1] \to [0,1]$ dado por:
	$$ f(x) = \begin{cases} 2x & : 0 \le x \le 1/2 \\ 2-2x & : 1/2 \le x \le 1 \end{cases} $$
	Sabemos que $f([0, 1/2]) = [0,1]$ y $f([1/2, 1]) = [0,1]$.
	De manera más general, las iteraciones sobre los subintervalos diádicos abarcan todo el espacio:
	$$ f^n\left(\left[\frac{i}{2^n}, \frac{i+1}{2^n}\right]\right) = [0,1], \quad \text{para } i=0, \dots, 2^n-1, \forall n \in \mathbb{N} $$
	Sean $U$ y $V \subseteq [0,1]$ abiertos. Existe un $n_0 \in \mathbb{N}$ tal que podemos encajar un subintervalo diádico entero dentro de $U$:
	$$ \exists n_0 \in \mathbb{N} / \left[\frac{i-1}{2^{n_0}}, \frac{i}{2^{n_0}}\right] \subseteq U \text{ para algún } i \in \{1, 2, \dots, 2^{n_0}\} $$
	Aplicando $f^{n_0}$ a ambos lados, obtenemos:
	$$ f^{n_0}\left(\left[\frac{i-1}{2^{n_0}}, \frac{i}{2^{n_0}}\right]\right) \subseteq f^{n_0}(U) \implies [0,1] \subseteq f^{n_0}(U) \implies f^{n_0}(U) = [0,1] $$
	En consecuencia, $f^n(U) = [0,1]$ para todo $n \ge n_0$. Como $V$ es no vacío, $f^n(U) \cap V = [0,1] \cap V = V \neq \emptyset$. Esto demuestra que el mapa tienda es topológicamente mixing.
\end{ejemplo}

\begin{ejercicio}
	Sea $f: S^1 \to S^1$ dada por $f(z) = z^2$. Mostrar que $f$ es topológicamente mixing.
\end{ejercicio}
\begin{proof}[Solución]
	Si utilizamos la representación equivalente en el intervalo $[0,1)$ mediante la identificación del ángulo, la función se comporta como $x \mapsto 2x \pmod 1$.
	Sean $U, V \subseteq S^1$ dos abiertos no vacíos. Como $U$ es un conjunto abierto, este debe contener al menos un pequeño arco continuo de longitud $\epsilon > 0$.
	Cada vez que aplicamos $f$, el ángulo se duplica, por lo que la longitud del arco resultante se multiplica por 2 (considerando solapamientos cuando pasa de $2\pi$). 
	Tras $n$ iteraciones, la imagen $f^n(U)$ abarcará una longitud proyectada de $2^n \epsilon$. 
	Dado que la circunferencia tiene una longitud finita de $2\pi$, podemos elegir un $n_0 \in \mathbb{N}$ lo suficientemente grande tal que:
	$$ 2^{n_0} \epsilon > 2\pi $$
	A partir de este $n_0$, el arco se ha expandido tanto que se enrolla cubriendo todo el círculo, es decir, $f^{n_0}(U) = S^1$. Evidentemente, esta cobertura total se mantendrá para todo $n \ge n_0$, por lo que $f^n(U) = S^1$.
	Finalmente, como $V$ no es vacío, deducimos que:
	$$ f^n(U) \cap V = S^1 \cap V = V \neq \emptyset \quad \forall n \ge n_0 $$
	Concluyendo así que $f(z) = z^2$ es topológicamente mixing.
\end{proof}

\begin{ejercicio}
	Sean $X, Y$ espacios topológicos y $f: X \to X$, $g: Y \to Y$ aplicaciones continuas tales que $f$ y $g$ son topológicamente conjugados vía $h: X \to Y$. Muestre que:
	\begin{enumerate}[1)]
		\item $f$ es topológicamente transitivo $\iff g$ es topológicamente transitivo.
		\item $f$ es topológicamente mixing $\iff g$ es topológicamente mixing.
	\end{enumerate}
\end{ejercicio}
\begin{proof}[Solución]
	Demostraremos la ida ($\Rightarrow$). La vuelta ($\Leftarrow$) es análoga aplicando las mismas propiedades al homeomorfismo inverso $h^{-1}$.
	
	\textbf{1)} Sean $U, V$ abiertos no vacíos en $Y$. Como $h$ es un homeomorfismo, sus preimágenes $h^{-1}(U)$ y $h^{-1}(V)$ son abiertos no vacíos en $X$.
	Dado que $f$ es topológicamente transitiva, existe $n \in \mathbb{N}$ tal que:
	$$ f^n(h^{-1}(U)) \cap h^{-1}(V) \neq \emptyset $$
	Aplicando la función $h$ a toda la intersección (y dado que al ser inyectiva preserva las intersecciones):
	$$ h\left(f^n(h^{-1}(U))\right) \cap h(h^{-1}(V)) \neq \emptyset $$
	Por la definición de conjugación, sabemos que $h \circ f^n = g^n \circ h$. Sustituyendo esto, obtenemos:
	$$ g^n(h(h^{-1}(U))) \cap V \neq \emptyset \implies g^n(U) \cap V \neq \emptyset $$
	Por ende, $g$ también es topológicamente transitiva.
	
	\textbf{2)} Nuevamente, sean $U, V$ abiertos no vacíos en $Y \implies h^{-1}(U), h^{-1}(V)$ son abiertos no vacíos en $X$.
	Como $f$ es topológicamente mixing, existe $n_0 \in \mathbb{N}$ tal que:
	$$ f^n(h^{-1}(U)) \cap h^{-1}(V) \neq \emptyset, \quad \forall n \ge n_0 $$
	Aplicando $h$ a la intersección bajo la misma lógica del punto anterior:
	$$ h\left(f^n(h^{-1}(U))\right) \cap h(h^{-1}(V)) \neq \emptyset, \quad \forall n \ge n_0 $$
	$$ g^n(U) \cap V \neq \emptyset, \quad \forall n \ge n_0 $$
	En consecuencia, $g$ es topológicamente mixing.
\end{proof}

\begin{ejercicio}
	Probar que para un espacio $X$ separable y localmente compacto, la rotación $R_\theta: S^1 \to S^1$ es topológicamente transitiva pero no topológicamente mixing (para $\theta \notin \mathbb{Q}$).
\end{ejercicio}
\begin{proof}[Solución]
	Ya demostramos en semanas anteriores que si el ángulo es irracional, la órbita es densa, garantizando que $R_\theta$ es topológicamente transitiva.
	
	Para probar que no es mixing, vayamos a la representación en el intervalo $[0,1)$, es decir, $R_\theta: [0,1) \to [0,1)$ con $x \mapsto R_\theta(x) = x + \theta \pmod 1$.
	Supongamos por contradicción que es topológicamente mixing $\implies \forall U, V$ abiertos en $[0,1)$, existe $n_0 \in \mathbb{N}$ tal que $R_\theta^n(U) \cap V \neq \emptyset$ para todo $n \ge n_0$.
	
	Vamos a dividir el intervalo $[0,1)$ en $32$ partes iguales. Sea $\varepsilon_0 = 1/32$. Tomemos estos abiertos específicos:
	$$ U = V_{\varepsilon_0}(1/2), \quad V = V_{\varepsilon_0}(1/2 + 12\varepsilon_0), \quad W = V_{\varepsilon_0}(1/2 - 12\varepsilon_0) $$
	Aplicando la suposición de mixing sobre los conjuntos $V$ y $U$, tenemos que $\exists n_0 \in \mathbb{N}$ tal que $R_\theta^{n_0}(V) \cap U \neq \emptyset$.
	
	Sea $\delta = \max\{|x-y| : x \in R_\theta^{n_0}(V) \land y \in U\}$. 
	Como $\exists z \in U \cap R_\theta^{n_0}(V)$, podemos utilizar la desigualdad triangular para acotar la máxima distancia entre cualquier elemento de $R_\theta^{n_0}(V)$ y $U$:
	$$ \delta \le \max\{|x-z| + |z-y| : x \in R_\theta^{n_0}(V), y \in U\} < 2\varepsilon_0 + 2\varepsilon_0 = 4\varepsilon_0 $$
	
	Ahora, tomemos un punto $x \in R_\theta^{n_0}(V)$ cualquiera. Utilizando nuevamente la suposición de mixing pero ahora cruzando $R_\theta^{n_0}(V)$ con $W$, sabemos que $\exists m_0 \in \mathbb{N}$ tal que $R_\theta^{m_0}(x) \in W$. 
	Por ende, $R_\theta^{n_0+m_0}(V) \cap W \neq \emptyset$. De forma análoga a la acotación anterior, la distancia máxima se preserva:
	$$ \max\{|x-y| : x \in R_\theta^{n_0+m_0}(V) \land y \in W\} < 4\varepsilon_0 $$
	
	Evaluemos el cruce $R_\theta^{m_0+n_0}(V) \cap U$, el cual, por la hipótesis de mixing (ya que $m_0+n_0 \ge n_0$), no debería ser vacío.
	Sea $x \in R_\theta^{m_0+n_0}(V)$. Por la acotación deducida con $W$, sabemos que $|x - (1/2 - 12\varepsilon_0)| < 4\varepsilon_0$.
	Si queremos que este mismo $x$ pertenezca también a la vecindad $U$, su distancia a $1/2$ tendría que cumplir $|x - 1/2| < 4\varepsilon_0$. Desarrollemos esta distancia:
	\begin{align*}
		|x - 1/2| &= |x - (1/2 - 12\varepsilon_0) - 12\varepsilon_0| \\
		&\ge |12\varepsilon_0| - |x - (1/2 - 12\varepsilon_0)| \\
		&> 12\varepsilon_0 - 4\varepsilon_0 = 8\varepsilon_0
	\end{align*}
	Llegamos a que $|x - 1/2| > 8\varepsilon_0$. Esto entra en conflicto directo con la condición obligatoria de que la distancia sea menor a $4\varepsilon_0$ ($\rightarrow \leftarrow$).
	Por tanto, tal intersección es imposible, es decir, $R_\theta^{m_0+n_0}(V) \cap U = \emptyset$. Esto rompe la definición general de iteración para todo $n \ge n_0$, demostrando que $R_\theta$ no es topológicamente mixing.
\end{proof}
	% =========================================================================
% --- INICIO INPUT SEMANA 3 ---
% =========================================================================

\chapter{Semana 3: Shift de Bernoulli y Semi-Conjugación}

\section{Shift de Bernoulli}

Sea $X = \{z_1, z_2, \dots, z_d\}$ un alfabeto finito equipado con la topología discreta generada por la métrica $d_1: X \times X \to \mathbb{R}^+_0$:
$$d_1(u, v) = \begin{cases} 0 & u = v \\ 1 & u \neq v \end{cases}$$
Definimos los espacios de sucesiones:
\begin{itemize}
	\item El espacio unilateral: $\Sigma^+ = X^{\mathbb{N}} = \{(x_n)_{n \in \mathbb{N}} : x_n \in X\}$
	\item El espacio bilateral: $\Sigma = X^{\mathbb{Z}} = \{(x_n)_{n \in \mathbb{Z}} : x_n \in X\}$
\end{itemize}

Podemos dotar a estos espacios de la siguiente métrica (escrita para $\Sigma$, de forma análoga para $\Sigma^+$):
$$d((x_n)_{n \in \mathbb{Z}}, (y_n)_{n \in \mathbb{Z}}) = \sum_{n \in \mathbb{Z}} \frac{d_1(x_n, y_n)}{2^{|n|}}$$

\begin{proposicion}
	La función $d$ es una métrica en $\Sigma$ (y en $\Sigma^+$).
\end{proposicion}
\begin{proof}
	Claramente está bien definida (converge), pues:
	$$ \sum_{n \in \mathbb{Z}} \frac{d_1(x_n, y_n)}{2^{|n|}} = \sum_{n \in \mathbb{N}} \frac{d_1(x_n, y_n)}{2^n} + \frac{d_1(x_0, y_0)}{2^0} + \sum_{n \in \mathbb{N}} \frac{d_1(x_{-n}, y_{-n})}{2^{-n}} \le \sum_{n \in \mathbb{N}} \frac{1}{2^n} + 1 + \sum_{n \in \mathbb{N}} \frac{1}{2^n} < 1 + 1 + 1 = 3 $$
	Los resultados siguientes son obvios porque $d_1$ es una métrica:
	\begin{enumerate}[i)]
		\item Si $(x_n)_{n \in \mathbb{Z}} = (y_n)_{n \in \mathbb{Z}} \implies \forall n \in \mathbb{Z} \implies x_n = y_n, \forall n \in \mathbb{Z} \implies d_1(x_n, y_n) = 0 \implies d((x_n), (y_n)) = 0$.
		Si $d((x_n), (y_n)) = 0 \implies d_1(x_n, y_n) = 0, \forall n \in \mathbb{Z} \implies x_n = y_n, \forall n \in \mathbb{Z} \implies (x_n) = (y_n)$.
		\item $d((x_n), (y_n)) = \sum_{n \in \mathbb{Z}} \frac{d_1(x_n, y_n)}{2^{|n|}} = \sum_{n \in \mathbb{Z}} \frac{d_1(y_n, x_n)}{2^{|n|}} = d((y_n), (x_n))$.
		\item Sea $(x_n), (y_n), (z_n) \in \Sigma \implies d((x_n), (y_n)) = \sum_{n \in \mathbb{Z}} \frac{d_1(x_n, y_n)}{2^{|n|}} \le \sum_{n \in \mathbb{Z}} \frac{d_1(x_n, z_n) + d_1(z_n, y_n)}{2^{|n|}} = d((x_n), (z_n)) + d((z_n), (y_n))$.
	\end{enumerate}
\end{proof}

\begin{definicion}[Cilindros]
	Un \textbf{cilindro} de longitud $k$ es el conjunto de todas las sucesiones que tienen símbolos fijos en coordenadas específicas. Lo denotamos por:
	$$[n_1: c_1, \dots, n_k: c_k] = \{(x_n) \in \Sigma : x_{n_i} = c_i \text{ para } i=1,\dots,k\}$$
\end{definicion}

\begin{observacion}
	Sea $X = \{z_1, z_2, \dots, z_d\}$ un conjunto finito y $\Sigma = X^{\mathbb{Z}}$.
	\begin{enumerate}[a)]
		\item $\Sigma = \bigsqcup_{i=1}^d [j: z_i]$ con $j \in \mathbb{Z}$ arbitrario y fijo.
		\item $\Sigma = \bigsqcup_{i,m=1}^d [j: z_i, z_m]$ con $j \in \mathbb{Z}$ arbitrario y fijo.
	\end{enumerate}
	En $X$ estamos considerando la topología generada por $d_1$. Así, $X$ es un espacio topológico compacto (Ejercicio).
	Como $\Sigma = X^{\mathbb{Z}}$ (y $\Sigma^+ = X^{\mathbb{N}}$), podemos considerar la topología producto en $\Sigma$, lo cual garantiza por el Teorema de Tychonoff que $\Sigma$ es compacto.
\end{observacion}

\begin{ejercicio}
	Mostrar que la colección de cilindros en $\Sigma^+$ (o en $\Sigma$) es una base topológica para la topología en $\Sigma^+$ (o en $\Sigma$) generada por $d$.
\end{ejercicio}
\begin{proof}[Solución]
	Sea $B = B_\epsilon((x_n)_{n \in \mathbb{Z}})$ con $\epsilon > 0$ y $(y_n)_{n \in \mathbb{Z}} \in B$. Sea $\delta = \min\{d(x,y), \epsilon - d(x,y)\}$. 
	Como $\delta > 0$, dividimos con un número $K \in \mathbb{N}$ tal que $0 < \delta/K = \delta^* < 1$, donde $K \ge 3$.
	Como $\delta^* < 1 \implies \exists n_0 \in \mathbb{N} / \frac{1}{2^{n_0-1}} = \sum_{n=n_0+1}^\infty \frac{1}{2^{n-1}} < \delta^*$. 
	Luego, definimos el cilindro: $C = [-n_0 : y_{-n_0}, \dots, n_0 : y_{n_0}]$. Claramente $(y_n)_{n \in \mathbb{N}} \in C$.
	
	Sea $(z_n)_{n \in \mathbb{Z}} \in C \implies d((z_n), (y_n)) = \sum_{n \in \mathbb{Z}} \frac{d_1(z_n, y_n)}{2^{|n|}} + \sum_{n \in \mathbb{N}} \frac{d_1(z_{-n}, y_{-n})}{2^n} = \sum_{n=n_0+1}^\infty \frac{d_1(z_n, y_n)}{2^n} + \sum_{n=n_0+1}^\infty \frac{d_1(z_{-n}, y_{-n})}{2^n} \le 2 \sum_{n=n_0+1}^\infty \frac{1}{2^n} < 2\delta^* < \delta$.
	Luego, $(z_n)_{n \in \mathbb{Z}} \in B_\delta(y) \implies (z_n)_{n \in \mathbb{Z}} \in C \subseteq B$.
\end{proof}

\begin{figure}[htbp]
	\centering
	\begin{tikzpicture}[scale=1.5]
		% Bola grande B (alrededor de x)
		\draw[blue, thick, dashed] (0,0) circle (2cm);
		\node[blue, right] at (1.4, 1.4) {$B = B_\epsilon(x)$};
		\filldraw[black] (0,0) circle (1pt) node[below left] {$x$};
		
		% Bola pequeña B_delta (alrededor de y)
		\draw[purple, thick, dashed] (0.8,0.5) circle (0.7cm);
		\node[purple, right] at (1.2, 0.9) {$B_\delta(y)$};
		\filldraw[black] (0.8,0.5) circle (1pt) node[right] {$y$};
		
		% Cilindro C (alrededor de y)
		\draw[orange, thick] (0.4, 0.1) rectangle (1.2, 0.9);
		\node[orange, left] at (0.4, 0.5) {$C$};
	\end{tikzpicture}
	\caption{El cilindro $C$ está contenido en la vecindad $B_\delta(y)$, la cual a su vez está contenida en el abierto original $B$.}
\end{figure}

\begin{ejercicio}
	Mostrar que cualquier cilindro $[m: a_m, \dots, a_n]$ es un conjunto abierto en $\Sigma^+$ (o en $\Sigma$) con la topología en $\Sigma^+$ (o en $\Sigma$) generada por $d$.
\end{ejercicio}
\begin{proof}[Solución]
	Sea $C = [m: a_m, \dots, a_n]$ un cilindro. Sea $x = (x_k) \in C$. 
	Queremos encontrar un $\epsilon > 0$ tal que $B_\epsilon(x) \subseteq C$. 
	Basta tomar $\epsilon < \frac{1}{2^{\max(|m|, |n|)}}$. Sea $y = (y_k) \in B_\epsilon(x)$, entonces $d(x, y) < \epsilon$.
	Si existiera algún $k \in [m, n]$ tal que $y_k \neq x_k = a_k$, entonces $d_1(x_k, y_k) = 1$.
	Esto implicaría que $d(x, y) \ge \frac{d_1(x_k, y_k)}{2^{|k|}} = \frac{1}{2^{|k|}} \ge \frac{1}{2^{\max(|m|, |n|)}} > \epsilon$, lo cual es una contradicción ($\rightarrow \leftarrow$).
	Por ende, $y_k = a_k$ para todo $k \in [m, n]$, lo que implica que $y \in C$. Así, $B_\epsilon(x) \subseteq C$, demostrando que el cilindro es un conjunto abierto.
\end{proof}

\begin{ejercicio}
	Mostrar que la topología en $\Sigma^+$ (en $\Sigma$) generada por $d$ es igual a la topología producto $\tau$.
\end{ejercicio}
\begin{proof}[Solución]
	i) Sea $U$ abierto en la topología producto: $U = B_{\epsilon_{i_1}} \times \dots \times B_{\epsilon_{i_n}} \times \prod_{j \neq i_1,\dots,i_n} X$, donde $\epsilon_{i_1}, \dots, \epsilon_{i_n} > 0$. Donde $j \in \mathbb{Z}$.
	Donde $x = (x_n)_{n \in \mathbb{Z}} \in U$. Definimos $C = [i_1: x_{i_1}, \dots, i_n: x_{i_n}]$. Luego, $x \in C \subseteq U$.
	
	ii) La forma análoga es muy sencilla. 
	Como la topología generada por $d$ es equivalente a la que es generada por cilindros, entonces es equivalente a la topología producto.
\end{proof}

\begin{definicion}[Shift de Bernoulli]
	Sea $X = \{z_1, \dots, z_d\}$, $\Sigma^+ = \{(x_n)_{n \in \mathbb{N}} : x_n \in X\}$ y $\Sigma = \{(x_n)_{n \in \mathbb{Z}} : x_n \in X\}$. Definimos:
	\begin{itemize}
		\item \textbf{Shift Unilateral:} $\sigma: \Sigma^+ \to \Sigma^+$ dado por $(x_n)_{n \in \mathbb{N}} \mapsto \sigma((x_n)_{n \in \mathbb{N}}) = (y_n)_{n \in \mathbb{N}}$, donde $y_n = x_{n+1}, \forall n \in \mathbb{N}$.
		\item \textbf{Shift Bilateral:} $\sigma: \Sigma \to \Sigma$ dado por $(x_n)_{n \in \mathbb{Z}} \mapsto \sigma((x_n)_{n \in \mathbb{Z}}) = (y_n)_{n \in \mathbb{Z}}$, donde $y_n = x_{n+1}, \forall n \in \mathbb{Z}$.
	\end{itemize}
\end{definicion}

\begin{ejemplo}
	Para el shift unilateral: Sea $X=\{0,1,2,3\}$, $\Sigma^+ = X^{\mathbb{N}}$.
	Si $(x_n)_{n \in \mathbb{N}} = (3,1,0,2,3,1,0,2,3,1,0,2,\dots) \implies (y_n)_{n \in \mathbb{N}} = \sigma((x_n)_{n \in \mathbb{N}}) = (1,0,2,3,1,0,2,3,1,0,2,\dots)$.
\end{ejemplo}

\begin{proposicion}
	El shift unilateral $\sigma: \Sigma^+ \to \Sigma^+$ cumple lo siguiente:
	\begin{enumerate}
		\item \textbf{Es continua.} En efecto, como los cilindros son los abiertos básicos en $\Sigma^+$, entonces es suficiente mostrar que $\sigma^{-1}([m: a_m, \dots, a_n])$ sea abierto en $\Sigma^+$.
		Observemos que $\sigma^{-1}([m: a_m, \dots, a_n]) = [m+1: a_m, \dots, a_n]$ pues, si $(z_n)_{n \in \mathbb{N}} \in \sigma^{-1}([m: a_m, \dots, a_n])$ entonces $(y_n) = \sigma((z_n)_{n \in \mathbb{N}}) \in [m: a_m, \dots, a_n]$.
		$y_m = a_m, y_{m+1} = a_{m+1}, \dots, y_n = a_n$.
		Como $y_k = z_{k+1} \implies z_{m+1} = a_m, z_{m+2} = a_{m+1}, \dots, z_{n+1} = a_n \implies ((z_n))_{n \in \mathbb{N}} \in [m+1: a_m, \dots, a_n]$.
		\item \textbf{No es inyectivo.} En efecto, $X = \{z_1, \dots, z_d\}$ y $\Sigma^+ = \{(x_n)_{n \in \mathbb{N}} : x_n \in X\}$. Sean $(x_n) = (z_1, z_1, z_2, z_2, z_2, \dots)$ y $(y_n) = (z_2, z_1, z_2, z_2, z_2, \dots)$.
		Vemos que $(x_n) \neq (y_n)$, pero $\sigma((x_n)) = (z_1, z_2, z_2, \dots) = \sigma((y_n))$.
		\item \textbf{Es sobre.} Para cualquier $(y_n)_{n \in \mathbb{N}} \in \Sigma^+$, consideremos $(x_n) = (z_1, y_0, y_1, \dots)$. Evidentemente $\sigma((x_n)) = (y_n)$, garantizando la sobreyectividad.
	\end{enumerate}
\end{proposicion}

\begin{ejercicio}
	Mostrar que el shift bilateral $\sigma: \Sigma \to \Sigma$ es un homeomorfismo.
\end{ejercicio}
\begin{proof}[Solución]
	\begin{itemize}
		\item \textbf{Inyectiva:} Sean $x = (x_n)_{n \in \mathbb{Z}}, y = (y_n)_{n \in \mathbb{Z}} \in \Sigma / \sigma(x) = \sigma(y) \implies (a_n)_{n \in \mathbb{Z}} = (b_n)_{n \in \mathbb{Z}}$ donde $a_n = x_{n+1} \land b_n = y_{n+1}, \forall n \in \mathbb{Z} \implies x_{n+1} = y_{n+1}, \forall n \in \mathbb{Z} \implies x = y$.
		\item \textbf{Sobre:} Sea $y = (y_n)_{n \in \mathbb{Z}} \implies \exists (x_n)_{n \in \mathbb{Z}} / x_n = y_{n-1}, \forall n \in \mathbb{Z} \implies \sigma((x_n)_{n \in \mathbb{Z}}) = (z_n)_{n \in \mathbb{Z}} / z_n = x_{n+1} = y_n, \forall n \in \mathbb{Z} \implies \sigma((x_n)) = (y_n)_{n \in \mathbb{Z}}$.
		(La inversa se define como $\sigma^{-1}: \Sigma \to \Sigma, \sigma^{-1}((x_n)_{n \in \mathbb{Z}}) = (y_n)_{n \in \mathbb{Z}} / y_n = x_{n-1}, \forall n \in \mathbb{Z}$).
		\item \textbf{Continua:} $\sigma: \Sigma \to \Sigma, \sigma^{-1}: \Sigma \to \Sigma$. Sea $C = [m: a_m, \dots, a_n]$ cualquier cilindro. $\sigma^{-1}(C) = [m+1: a_m, \dots, a_n] \land (\sigma^{-1})^{-1}(C) = \sigma(C) = [m-1: a_m, \dots, a_n]$ cilindros.
	\end{itemize}
\end{proof}

\begin{proposicion}
	Para el shift bilateral, $\overline{Per(\sigma)} = \Sigma$.
\end{proposicion}
\begin{proof}
	Como los abiertos básicos de $\Sigma$ son los cilindros, entonces es suficiente mostrar que $Per(\sigma) \cap [m: a_m, \dots, a_n] \neq \emptyset$.
	Consideremos: $(x_n)_{n \in \mathbb{Z}} = (\dots a_m \dots a_n \dots a_m \dots a_n \dots) \in [m: a_m, \dots, a_n]$ (bloque central de tamaño $n-m+1$) y evidentemente $\sigma^{n-m+1}((x_n)_{n \in \mathbb{Z}}) = (x_n)_{n \in \mathbb{Z}}$.
\end{proof}

\begin{ejercicio}
	¿Se cumple $\overline{Per(\sigma^+)} = \Sigma^+$?
\end{ejercicio}
\begin{proof}[Solución]
	Sea $x = (x_n)_{n \in \mathbb{N}} \in \Sigma^+ \land C = [m: x_m, \dots, x_r]$ con $m \ge 1$. Definimos: $y = (y_n)_{n \in \mathbb{N}} / y_{n+r} = y_n, \forall n \in \mathbb{N}$. Luego: $\sigma^r(y) = y \implies y \in Per(\sigma^+)$. Además $y \in C$, con $y_m = x_m, \dots, y_r = x_r$. Luego: $\overline{Per(\sigma^+)} = \Sigma^+$.
\end{proof}

\begin{proposicion}
	El shift bilateral $\sigma: \Sigma \to \Sigma$ es topológicamente mixing.
\end{proposicion}
\begin{proof}
	Sean $U, V \subseteq \Sigma$ abiertos en $\Sigma$ como los cilindros son los abiertos básicos de $\Sigma$, entonces existen $[p: a_p, \dots, a_q] \subseteq U, [r: b_r, \dots, b_s] \subseteq V$.
	Afirmación: $\exists n_0 \in \mathbb{N} / \sigma^n([p: a_p, \dots, a_q]) \cap [r: b_r, \dots, b_s] \neq \emptyset, \forall n \ge n_0$. Así $\sigma^n([p: a_p, \dots, a_q]) \cap [r: b_r, \dots, b_s] \subseteq \sigma^n(U) \cap V \neq \emptyset, \forall n \ge n_0$.
	Prueba de la afirmación: Existe $m \in \mathbb{N}$ y $\alpha_{-m}, \dots, \alpha_m, \beta_{-m}, \dots, \beta_m \in X$ tal que:
	$$ C_1^m = [-m: \alpha_{-m}, \dots, \alpha_m] \subseteq [p: a_p, \dots, a_q] $$
	$$ C_2^m = [-m: \beta_{-m}, \dots, \beta_m] \subseteq [r: b_r, \dots, b_s] $$
	¡Ejercicio! (Basta tomar $m = \max\{|p|, |q|, |r|, |s|\}$).
	
	Afirmación 1: Tomando $n_0 = 2m+2$. Se tiene que $\sigma^n(C_1^m) \cap C_2^m \neq \emptyset, \forall n \ge n_0$. En efecto, para cada $n = 2m+1+k, k \in \mathbb{Z}^+$ definimos $(w_i)_{i \in \mathbb{Z}}$:
	$w_i = \alpha_i$ para $|i| \le m$, y $w_i = \beta_{i-n}$ para $i = m+k+1, \dots, 3m+k+1$.
	Luego: $(w_i)_{i \in \mathbb{Z}} \in C_1^m \implies \sigma^{2m+1+k}((w_i)_{i \in \mathbb{Z}}) \in C_2^m \implies \sigma^n(C_1^m) \cap C_2^m \neq \emptyset$. Como $k$ fue arbitrario en $\mathbb{Z}^+ \implies \sigma^n(C_1^m) \cap C_2^m \neq \emptyset, \forall n \ge n_0 = 2m+2$.
\end{proof}

\begin{corolario}
	$\sigma: \Sigma \to \Sigma$ es topológicamente transitivo. (Pues $\Sigma$ es métrico, compacto separable y localmente compacto).
\end{corolario}

\section{Semi-Conjugación Topológica}

\begin{definicion}[Semi-Conjugación]
	Sean $X, Y$ espacios topológicos y $f: X \to X$, $g: Y \to Y$ aplicaciones continuas. Decimos que $g$ es un \textbf{factor} de $f$ (o que $f$ es semiconjugada a $g$) si existe una función $h: X \to Y$ \textbf{continua y sobreyectiva} tal que:
	$$h \circ f = g \circ h$$
\end{definicion}

\begin{ejemplo}
	Consideremos el shift unilateral $\Sigma^+$ en el alfabeto $X=\{0,1\}$, y la aplicación de ángulo doble $f(x) = 2x \pmod 1$ en el espacio $[0,1]$.
	Definimos la función $h: \Sigma^+ \to [0,1]$ como la expansión binaria:
	$$h((x_n)) = \sum_{n=0}^{\infty} \frac{x_n}{2^{n+1}}$$
	Esta función $h$ es una semi-conjugación. Es sobreyectiva porque todo número real en $[0,1]$ tiene al menos una representación binaria (ilustrado en la Figura \ref{fig:expansion_binaria}). No es inyectiva porque los números diádicos tienen dos representaciones (por ejemplo, $0.0111\dots = 0.1000\dots = 1/2$).
	
	\begin{figure}[htbp]
		\centering
		\begin{tikzpicture}[scale=4]
			% Eje principal
			\draw[thick] (0,0) -- (1,0);
			
			% Nivel 0 (0 y 1)
			\draw[thick] (0,-0.1) node[below] {$0$} -- (0,0.1);
			\draw[thick] (1,-0.1) node[below] {$1$} -- (1,0.1);
			
			% Nivel 1 (Mitades)
			\draw[thick, blue] (0.5,-0.1) node[below, text=black] {$1/2$} -- (0.5,0.1);
			\node[above, blue] at (0.25, 0.1) {$.0$};
			\node[above, blue] at (0.75, 0.1) {$.1$};
			
			% Nivel 2 (Cuartos)
			\draw[thick, red] (0.25,-0.05) node[below, text=black, yshift=-0.5cm] {$1/4$} -- (0.25,0.05);
			\draw[thick, red] (0.75,-0.05) node[below, text=black, yshift=-0.5cm] {$3/4$} -- (0.75,0.05);
			\node[above, red] at (0.125, 0.3) {$.00$};
			\node[above, red] at (0.375, 0.3) {$.01$};
			\node[above, red] at (0.625, 0.3) {$.10$};
			\node[above, red] at (0.875, 0.3) {$.11$};
			
			% Conectores para mostrar pertenencia a subintervalos
			\draw[->, gray, thin] (0.125, 0.25) -- (0.125, 0.05);
			\draw[->, gray, thin] (0.375, 0.25) -- (0.375, 0.05);
			\draw[->, gray, thin] (0.625, 0.25) -- (0.625, 0.05);
			\draw[->, gray, thin] (0.875, 0.25) -- (0.875, 0.05);
		\end{tikzpicture}
		\caption{Interpretación de $h((x_n))$ como expansión binaria. Los dígitos de la sucesión $(x_0, x_1, \dots)$ especifican sucesivamente en qué mitad del intervalo cae el punto.}
		\label{fig:expansion_binaria}
	\end{figure}
	
	Comprobamos la condición de semiconjugación $h \circ \sigma = f \circ h$:
	\begin{align*}
		h(\sigma(x_0, x_1, x_2 \dots)) &= h(x_1, x_2, \dots) = \sum_{n=0}^{\infty} \frac{x_{n+1}}{2^{n+1}} \\
		f(h(x_0, x_1, \dots)) &= 2 \left( \sum_{n=0}^{\infty} \frac{x_n}{2^{n+1}} \right) \pmod 1 \\
		&= \left( x_0 + \sum_{n=1}^{\infty} \frac{x_n}{2^n} \right) \pmod 1
	\end{align*}
	Como $x_0 \in \{0,1\}$, al aplicar módulo 1 la parte entera desaparece, quedando exactamente la misma suma que obtuvimos para $h(\sigma(x))$. Así, la dinámica del shift sobre símbolos emula la dinámica de $2x \pmod 1$.
\end{ejemplo}

\begin{ejercicio}
	Muestre que si $f$ es topológicamente transitiva y $g$ es un factor de $f$ (es decir, existe una semiconjugación $h$), entonces $g$ también es topológicamente transitiva.
\end{ejercicio}
\begin{proof}[Solución]
	Sea $h: X \to Y$ la semiconjugación sobreyectiva y continua ($h \circ f = g \circ h$). Como $f$ es transitiva, existe un punto $x \in X$ con órbita densa, es decir $\overline{\mathcal{O}_f^+(x)} = X$.
	Queremos evaluar la órbita del punto $y = h(x)$ bajo la función $g$.
	Aplicando $h$ a la órbita densa, por continuidad de $h$ tenemos:
	$$h\left(\overline{\mathcal{O}_f^+(x)}\right) \subseteq \overline{h(\mathcal{O}_f^+(x))}$$
	Reemplazando la clausura en la izquierda, $h(X) \subseteq \overline{h(\mathcal{O}_f^+(x))}$. Por ser $h$ sobreyectiva, $h(X) = Y$.
	Además, evaluemos los puntos de la órbita mapeada:
	$h(f^n(x)) = g^n(h(x)) = g^n(y)$.
	Sustituyendo esto, obtenemos que $Y \subseteq \overline{\{g^n(y)\}_{n \ge 0}} = \overline{\mathcal{O}_g^+(y)}$.
	Como el espacio entero $Y$ está contenido en la clausura de la órbita de $y$, dicha órbita es densa, lo que demuestra que $g$ es topológicamente transitiva.
\end{proof}

% =========================================================================
% --- FIN INPUT SEMANA 3 ---
% =========================================================================
	

	
\end{document}